\section{AI-Driven Dynamic Narrative Personalization in Games: A
Closed-Loop Plugin Using Hexad Player Modeling, RAG-Grounded LLM
Generation, and PPO Reinforcement
Learning}\label{ai-driven-dynamic-narrative-personalization-in-games-a-closed-loop-plugin-using-hexad-player-modeling-rag-grounded-llm-generation-and-ppo-reinforcement-learning}

\begin{center}\rule{0.5\linewidth}{0.5pt}\end{center}

\textbf{Candidate:} {[}Candidate Name{]} \textbf{Supervisor:}
{[}Supervisor Name{]} \textbf{Institution:} {[}University Name{]}
\textbf{Department:} Computer Science / Human-Computer Interaction
\textbf{Degree:} Master of Science (MSc) \textbf{Submission Date:}
February 2026 \textbf{Word Count:} {[}approx. 12,500 words{]}

\begin{center}\rule{0.5\linewidth}{0.5pt}\end{center}

\emph{This dissertation is submitted in partial fulfilment of the
requirements for the degree of Master of Science. The work presented is
entirely the candidate's own, except where otherwise stated.}

\begin{center}\rule{0.5\linewidth}{0.5pt}\end{center}

\subsection{Abstract}\label{abstract}

Narrative personalization in games --- the capacity of a game's story to
adapt meaningfully to individual player behaviour --- remains an
under-explored intersection of artificial intelligence, human-computer
interaction, and game design. Most commercial games deliver identical
narrative content regardless of how a player engages with the world,
foregoing the well-established link between personalized experience and
intrinsic motivation.

This dissertation presents the design, implementation, and evaluation of
a closed-loop Unity plugin for real-time dynamic narrative
personalization. The system integrates five tightly coupled modules: a
telemetry collector that streams 23 behavioural signals from the game
client every five seconds; a player modeling service that derives an
11-feature representation, a six-dimensional continuous Hexad player
type profile, and a Flow state classification (FLOW, BOREDOM, ANXIETY,
APATHY) from those signals; a narrative generation service that employs
Retrieval-Augmented Generation (RAG) over a structured lore corpus using
the sentence-transformers all-MiniLM-L6-v2 model together with an
episodic session memory store; an adaptation engine based on Proximal
Policy Optimization (PPO), trained via Stable-Baselines3 over a custom
Gymnasium Markov Decision Process environment; and a Unity Content
Injection Module that surfaces the generated narrative inside a 2D
top-down RPG testbed set in the world of Eryndal.

A key novelty is the derivation of Hexad player type scores entirely
from gameplay telemetry --- without requiring the player to complete a
questionnaire --- enabling continuous, within-session profiling. The PPO
agent selects among eight narrative actions (e.g., PROVIDE\_GUIDANCE,
ADD\_MYSTERY, INCREASE\_URGENCY) based on the current Flow state and
Hexad profile, with a reward function shaped around challenge-skill
balance. A cold-start heuristic governs the first 120 seconds before the
learned policy activates.

The system is evaluated in a between-subjects user study (N = 50; 25
Experimental, 25 Control). The primary dependent variable is the Flow
subscale of the Game Experience Questionnaire (GEQ); secondary measures
include the GEQ Immersion subscale, the custom three-item Narrative
Personalization Scale (NPS-3), and the miniPXI player experience
inventory. Qualitative data are gathered through semi-structured
interviews analyzed using Reflexive Thematic Analysis (Braun \& Clarke,
2022). Results are pending collection and are presented as structured
placeholders.

The work makes six original contributions: telemetry-derived Hexad
profiling, RAG-grounded narrative generation, episodic session memory in
the LLM prompt, PPO for narrative action selection, a Flow-based MDP
reward, and a complete closed-loop integration of all five modules
within a single game plugin.

\textbf{Keywords:} narrative personalization, player modeling, Hexad
typology, Flow theory, reinforcement learning, large language models,
retrieval-augmented generation, game experience, Unity plugin.

\begin{center}\rule{0.5\linewidth}{0.5pt}\end{center}

\subsection{Table of Contents}\label{table-of-contents}

\begin{itemize}
\tightlist
\item
  \hyperref[abstract]{Abstract}
\item
  \hyperref[chapter-1-introduction]{Chapter 1: Introduction}

  \begin{itemize}
  \tightlist
  \item
    1.1 Background and Motivation
  \item
    1.2 Problem Statement
  \item
    1.3 Research Questions
  \item
    1.4 Aims and Objectives
  \item
    1.5 Scope and Limitations
  \item
    1.6 Thesis Structure
  \end{itemize}
\item
  \hyperref[chapter-2-literature-review]{Chapter 2: Literature Review}

  \begin{itemize}
  \tightlist
  \item
    2.1 Narrative Personalization in Games
  \item
    2.2 Player Modeling Approaches
  \item
    2.3 Flow Theory in Game Design
  \item
    2.4 Large Language Models for Game Narrative
  \item
    2.5 Reinforcement Learning for Game Adaptation
  \item
    2.6 Episodic Memory for LLMs
  \item
    2.7 Research Gap and Positioning
  \end{itemize}
\item
  \hyperref[chapter-3-research-methodology]{Chapter 3: Research
  Methodology}

  \begin{itemize}
  \tightlist
  \item
    3.1 Research Design Overview
  \item
    3.2 System Architecture
  \item
    3.3 Player Modeling Module
  \item
    3.4 Narrative Generation Module
  \item
    3.5 Adaptation Engine
  \item
    3.6 Unity Plugin Design
  \item
    3.7 Testbed Game Design
  \item
    3.8 Evaluation Design
  \end{itemize}
\item
  \hyperref[chapter-4-system-implementation]{Chapter 4: System
  Implementation}

  \begin{itemize}
  \tightlist
  \item
    4.1 Technology Stack
  \item
    4.2 Backend Service
  \item
    4.3 Database Layer
  \item
    4.4 Player Modeling Implementation
  \item
    4.5 Narrative Pipeline Implementation
  \item
    4.6 RL Adaptation Implementation
  \item
    4.7 Unity Plugin Implementation
  \item
    4.8 Testing and Validation
  \end{itemize}
\item
  \hyperref[chapter-5-results]{Chapter 5: Results}

  \begin{itemize}
  \tightlist
  \item
    5.1 Participant Demographics
  \item
    5.2 Quantitative Results
  \item
    5.3 Qualitative Results
  \item
    5.4 Summary of Findings
  \end{itemize}
\item
  \hyperref[chapter-6-discussion]{Chapter 6: Discussion}

  \begin{itemize}
  \tightlist
  \item
    6.1 Interpretation of Quantitative Findings
  \item
    6.2 Themes from Interviews
  \item
    6.3 Novelty and Contributions
  \item
    6.4 Limitations
  \item
    6.5 Threats to Validity
  \end{itemize}
\item
  \hyperref[chapter-7-conclusion]{Chapter 7: Conclusion}

  \begin{itemize}
  \tightlist
  \item
    7.1 Summary of Contributions
  \item
    7.2 Future Work
  \item
    7.3 Closing Remarks
  \end{itemize}
\item
  \hyperref[references]{References}
\item
  \hyperref[appendices]{Appendices}
\end{itemize}

\begin{center}\rule{0.5\linewidth}{0.5pt}\end{center}

\subsection{List of Figures}\label{list-of-figures}

\begin{itemize}
\tightlist
\item
  \textbf{Figure 1} --- High-level system architecture (five-module
  closed loop)
\item
  \textbf{Figure 2} --- Player modeling data flow: telemetry to Hexad
  profile and Flow state
\item
  \textbf{Figure 3} --- Narrative generation pipeline: RAG retrieval,
  episodic memory, and LLM generation
\item
  \textbf{Figure 4} --- Markov Decision Process (MDP) state diagram for
  the adaptation engine
\item
  \textbf{Figure 5} --- Cold-start heuristic decision tree
\item
  \textbf{Figure 6} --- Unity plugin component diagram
  (PlayerDataLogger, NarrativeManager, ContentInjector)
\item
  \textbf{Figure 7} --- Testbed game map: Ashfen Village Hub, Dungeon
  Room A, Dungeon Room B
\item
  \textbf{Figure 8} --- Study procedure timeline per participant
\end{itemize}

\begin{center}\rule{0.5\linewidth}{0.5pt}\end{center}

\subsection{List of Tables}\label{list-of-tables}

\begin{itemize}
\tightlist
\item
  \textbf{Table 1} --- Summary of related work on narrative
  personalization and player modeling
\item
  \textbf{Table 2} --- TelemetryBatch schema (23 fields)
\item
  \textbf{Table 3} --- Eleven extracted features and computation
  formulae
\item
  \textbf{Table 4} --- Hexad dimension computation weights from
  telemetry signals
\item
  \textbf{Table 5} --- Flow state classification decision rules
\item
  \textbf{Table 6} --- Eight narrative actions with descriptions and
  intended effects
\item
  \textbf{Table 7} --- MDP observation vector (7 dimensions)
\item
  \textbf{Table 8} --- Reward function components and values
\item
  \textbf{Table 9} --- PPO hyperparameters
\item
  \textbf{Table 10} --- GEQ subscale structure (items and scoring)
\item
  \textbf{Table 11} --- Participant demographics {[}INSERT
  placeholder{]}
\item
  \textbf{Table 12} --- Descriptive statistics and t-test results for
  all DVs {[}INSERT placeholder{]}
\item
  \textbf{Table 13} --- Qualitative themes and representative quotes
  {[}INSERT placeholder{]}
\end{itemize}

\begin{center}\rule{0.5\linewidth}{0.5pt}\end{center}

\subsection{Chapter 1: Introduction}\label{chapter-1-introduction}

\subsubsection{1.1 Background and
Motivation}\label{background-and-motivation}

Games occupy a unique position in interactive media: they are the only
narrative form in which the audience simultaneously controls the
protagonist and shapes the story's unfolding. Yet despite the enormous
expressive potential this interactivity affords, the vast majority of
commercially released titles deliver largely identical narrative
sequences to every player. The village elder delivers the same warning,
the antagonist delivers the same monologue, and the quest log records
the same objectives regardless of whether the player is a cautious
explorer who reads every piece of environmental lore or an aggressive
disruptor who skips dialogue and charges headlong into combat.

This homogeneity is not a result of disinterest from game designers ---
narrative branching is a well-recognized craft challenge --- but rather
a structural consequence of the cost of authoring alternatives. Every
branch, every personalized response, requires writing, voice-acting, and
integration effort that scales combinatorially. The practical ceiling
for authored branching in even the most narrative-ambitious titles
(e.g., CD Projekt Red's Cyberpunk 2077 or Larian Studios' Baldur's Gate
3) is reached long before true player-adaptive personalization becomes
possible.

The emergence of large language models (LLMs) capable of coherent
long-form natural language generation, combined with advances in
real-time player behavior modeling and reinforcement learning-based
policy optimization, suggests a new possibility: narrative that is
generated, rather than authored, in response to the individual player's
demonstrated preferences, cognitive state, and engagement patterns. If
such a system can be made reliable, grounded in lore-consistent content,
and responsive to psychological signals of engagement, it could
represent a qualitative shift in the nature of player experience.

This dissertation takes up that challenge. It presents a closed-loop
plugin --- implementable within Unity --- that continuously profiles the
player using a telemetry-derived Hexad player type model (Tondello et
al., 2016) and a Flow state classifier grounded in Csikszentmihalyi's
(1990) challenge-skill balance theory. A Proximal Policy Optimization
(PPO; Schulman et al., 2017) agent selects among eight narrative
modulation actions based on these real-time signals. A
Retrieval-Augmented Generation (RAG) pipeline grounds the LLM's output
in a structured lore corpus, and an episodic session memory component
inspired by EM-LLM (Fountas et al., 2024) ensures narrative continuity
across a session.

The resulting system is evaluated against a static narrative control in
a between-subjects user study, measuring Flow experience, immersion, and
perception of personalization using established psychometric
instruments.

\subsubsection{1.2 Problem Statement}\label{problem-statement}

The central problem this research addresses is the following:
\textbf{current game narrative systems are static in the sense that they
do not adapt their tone, urgency, or content to the psychological state
and demonstrated preferences of the individual player in real time.}
This creates a mismatch between player engagement state and narrative
offering. A player experiencing BOREDOM --- whose challenge-skill ratio
has fallen below a meaningful threshold --- continues to receive the
same measured, low-stakes descriptions that a Flow-state player
receives. A player in ANXIETY --- overwhelmed by challenge --- receives
no additional guidance from the story world. A player whose behavior
reveals deep curiosity about lore receives the same minimal world
description as one who ignores it entirely.

The consequences of this mismatch are reduced engagement, diminished
immersion, and missed opportunities to use narrative as a lever for
returning players to a state of optimal experience. The problem is not
merely aesthetic: motivation research (Ryan \& Deci, 2000) demonstrates
that perceived competence, autonomy, and relatedness are foundational to
sustained intrinsic engagement, and narrative personalization is a
direct mechanism for signaling each of these.

The challenge of solving this problem is threefold. First, reliable
real-time player state inference requires robust signal processing from
noisy behavioral telemetry. Second, LLM-generated narrative in games
must be tightly constrained to avoid breaking fictional immersion
through anachronism, fourth-wall violations, or lore inconsistency.
Third, the adaptation policy must be learned rather than fully
hand-crafted, because the mapping from player state to optimal narrative
action is complex and context-dependent.

\subsubsection{1.3 Research Questions}\label{research-questions}

This dissertation is organized around three research questions:

\textbf{RQ1:} Can a closed-loop AI system driven by continuous Hexad
profiling and Flow-state detection produce narrative that is perceived
as more personalized than static pre-authored narrative?

\textbf{RQ2:} Does PPO-driven narrative adaptation result in measurably
higher Flow experience (as measured by the GEQ Flow subscale) compared
to a control condition receiving static narrative?

\textbf{RQ3:} How do players perceive and experience dynamically
personalized narrative in terms of engagement and immersion?

RQ1 addresses the subjective perception of personalization. RQ2 targets
the objective psychological measure of Flow, operationalized through a
validated instrument. RQ3 is exploratory and captures nuances that
quantitative measures may not surface.

\subsubsection{1.4 Aims and Objectives}\label{aims-and-objectives}

The overarching aim of this research is to demonstrate that a
closed-loop, AI-driven narrative personalization system embedded within
a game engine plugin can measurably improve player experience relative
to static narrative delivery.

The specific objectives are:

\begin{enumerate}
\def\labelenumi{\arabic{enumi}.}
\tightlist
\item
  To design and implement a five-module pipeline that integrates
  telemetry collection, player modeling, narrative generation, RL-based
  adaptation, and in-game content injection within a single coherent
  system.
\item
  To develop a method for deriving a six-dimensional Hexad player type
  profile from behavioral telemetry without requiring a player
  questionnaire.
\item
  To construct a Gymnasium-compatible MDP environment that models
  narrative adaptation as a sequential decision problem with a
  Flow-based reward signal.
\item
  To build a RAG pipeline grounded in a structured game lore corpus that
  constrains LLM output to in-world content.
\item
  To design and conduct a between-subjects user study (N = 50) using
  validated psychometric instruments to evaluate the system's impact on
  Flow, immersion, and perceived personalization.
\item
  To contribute qualitative insight through reflexive thematic analysis
  of post-play interviews.
\end{enumerate}

\subsubsection{1.5 Scope and Limitations}\label{scope-and-limitations}

The scope of this work is bounded in several important respects. The
testbed is a purpose-built 2D top-down RPG (set in the world of Eryndal)
rather than a commercially released game; this ensures experimental
control but limits ecological validity. The Hexad profiling method is
rule-based and heuristic rather than a trained classifier; this is a
deliberate design choice motivated by the cold-start constraint (no
pre-existing player data) but introduces approximation error relative to
ground-truth questionnaire-derived types.

The LLM employed is Phi-3.5 Mini (3.8B parameters), a locally hosted
model via Ollama. This choice prioritizes privacy and latency
predictability over raw generation quality; larger cloud-hosted models
would likely yield more fluent output but introduce dependency, cost,
and data governance concerns. The PPO agent is trained on a simulated
environment rather than directly on human play data, which means the
learned policy reflects the simulator's transition dynamics rather than
empirical player behavior.

The user study is bounded to N = 50, which is adequate for detecting
medium-to-large effect sizes (d ≥ 0.5) but would be insufficient to
detect subtle effects. The study uses a between-subjects design, which
eliminates carry-over effects but requires careful matching of
conditions for comparable baseline experience. Results from Chapter 5
onward are presented as structured templates pending data collection.

\subsubsection{1.6 Thesis Structure}\label{thesis-structure}

Chapter 2 reviews the relevant literature across narrative
personalization, player modeling, Flow theory, LLMs for games,
reinforcement learning for adaptation, and episodic memory. Chapter 3
presents the research methodology in detail, covering the system
architecture, each module's design, the testbed game, and the evaluation
design including hypotheses, instruments, procedure, and analysis plan.
Chapter 4 documents the technical implementation. Chapter 5 presents the
results of the user study (as structured placeholders). Chapter 6
discusses the findings, contributions, limitations, and threats to
validity. Chapter 7 concludes with a summary of contributions, future
work directions, and closing remarks. Full questionnaire instruments,
the API specification, and the MDP formal specification are provided in
the appendices.

\begin{center}\rule{0.5\linewidth}{0.5pt}\end{center}

\subsection{Chapter 2: Literature
Review}\label{chapter-2-literature-review}

\subsubsection{2.1 Narrative Personalization in
Games}\label{narrative-personalization-in-games}

The question of how games might deliver individualized narrative
experiences has been explored for several decades. Early work on
interactive narrative focused on structural approaches: drama management
systems such as the Facade project (Mateas \& Stern, 2003) attempted to
guide story structure toward dramatically satisfying outcomes by
monitoring player actions and intervening in the plot. These systems
were computationally tractable but brittle --- they required exhaustive
hand-authoring of alternative branches and could not generalize to novel
player behaviors.

More recent work has approached the problem from a content generation
perspective. Procedural narrative generation systems, such as Tracery
(Compton et al., 2015) and the Expressionist framework (Kreminski et
al., 2019), demonstrated that grammatical expansion could produce
contextually varied textual content. However, such systems still
required extensive authoring of production rules and lacked the capacity
to respond to the player's psychological state.

The commercial industry has converged on a middle path: narrative choice
systems with authored branches (e.g., Mass Effect, The Witcher series,
Disco Elysium) provide the illusion of personalization through
meaningful choices but are ultimately deterministic in their narrative
delivery relative to choice history. These systems do not adapt to
behavioral signals such as play pace, combat avoidance, or exploration
rate --- they adapt only to explicit player-authored choices within a
pre-specified decision space.

The emerging paradigm, exemplified by PANGeA (Todd et al., 2024) and
LIGS (see Section 2.4), uses generative AI to produce narrative content
at runtime, opening the possibility of adaptation that is both infinite
in variety and grounded in the game's content universe. The current work
situates itself at the intersection of this generative paradigm and
player modeling: the narrative is not merely generated but selected and
styled in response to a continuous model of the individual player's
state.

\subsubsection{2.2 Player Modeling
Approaches}\label{player-modeling-approaches}

\paragraph{2.2.1 Bartle Taxonomy and Its
Limitations}\label{bartle-taxonomy-and-its-limitations}

Bartle's (1996) taxonomy of MUD player types --- Achievers, Explorers,
Socializers, and Killers --- was the first systematic attempt to
categorize players by behavioral motivation. Developed through
observation of Multi-User Dungeon players in the 1980s, the taxonomy
introduced the foundational insight that players seek qualitatively
different experiences from the same game: some prioritize accumulation
and achievement, others exploration and discovery, others social
interaction, and others competitive dominance.

Despite its wide adoption in game design pedagogy, the Bartle taxonomy
has several well-documented limitations. First, it is categorical rather
than continuous: a player is assigned to one type, when empirical
evidence suggests players express multiple motivations simultaneously
and in varying proportions (Yee, 2006). Second, the taxonomy was derived
from a narrow population playing a specific genre (text-based MUDs) and
has questionable generalizability to modern games. Third, the categories
are not mutually exclusive and lack operationalized measurement
criteria. Finally, type assignment in the Bartle framework relies on
self-report questionnaires (the Bartle Test), which introduce social
desirability bias and are one-time snapshots rather than dynamic,
within-session measures.

\paragraph{2.2.2 Hexad Typology (Tondello et al.,
2016)}\label{hexad-typology-tondello-et-al.-2016}

Tondello et al.~(2016) developed the Gamification User Types Hexad Scale
as a theoretically grounded alternative to Bartle's taxonomy. Rooted in
Self-Determination Theory (Ryan \& Deci, 2000), the Hexad framework
identifies six player types: Achievers (motivated by mastery and
progress), Explorers (motivated by discovery and curiosity), Socializers
(motivated by social interaction and belonging), Free Spirits (motivated
by autonomy and self-expression), Disruptors (motivated by challenging
the status quo), and Philanthropists (motivated by giving and altruism).
Unlike Bartle's types, Hexad scores are continuous (typically measured
on a seven-point Likert scale per item) and can be meaningfully
combined.

The original Hexad scale is a 24-item self-report questionnaire,
validated for gamification contexts with acceptable reliability
(Cronbach's alpha typically exceeding 0.70 per subscale). Subsequent
work has applied Hexad profiling to adaptive game design, recommendation
systems, and educational technology, establishing it as the de facto
standard for research-grade player type assessment.

The critical limitation for the current work is the reliance on a
pre-play questionnaire. A questionnaire-derived profile is a static
snapshot, captured before any gameplay, that cannot update as the
player's engagement patterns evolve during a session. Moreover,
administering a 24-item questionnaire before play introduces friction
that reduces ecological validity. The present work addresses this
limitation by inferring Hexad-analogous continuous scores directly from
behavioral telemetry --- a novel contribution detailed in Section 3.3.

\paragraph{2.2.3 Data-Driven Player Modeling
(player2vec)}\label{data-driven-player-modeling-player2vec}

The player2vec line of research (represented by an arXiv preprint from
April 2024) demonstrates the application of Transformer-based
self-supervised learning to player behavior sequences. Analogous to
word2vec embeddings for natural language tokens, player2vec encodes
sequences of in-game actions into dense low-dimensional representations
that cluster players by behavioral similarity without requiring labeled
training data. This approach offers the compelling advantage of
discovering latent behavior patterns that pre-specified taxonomies may
miss.

The current work draws conceptual inspiration from player2vec ---
specifically, the notion that a session embedding can summarize a
player's behavioral trajectory --- but implements a simpler rule-based
approach rather than a trained Transformer. This pragmatic choice is
motivated by the cold-start constraint: a Transformer session embedder
requires a substantial corpus of historical play sequences to produce
meaningful representations, which is unavailable for a first-session
player using a newly deployed plugin.

\subsubsection{2.3 Flow Theory in Game Design (Csikszentmihalyi,
1990)}\label{flow-theory-in-game-design-csikszentmihalyi-1990}

Csikszentmihalyi's (1990) theory of Flow describes a psychological state
of optimal experience characterized by complete absorption in a
challenging activity, loss of self-consciousness, and distorted time
perception. Flow arises at the intersection of high challenge and high
skill: when challenge substantially exceeds skill, the result is
anxiety; when skill substantially exceeds challenge, the result is
boredom; and when both are absent, the result is apathy.

In game design, Flow theory has been influential since Chen's (2007)
canonical analysis of game flow, which argued that the dynamic
difficulty adjustment (DDA) mechanisms present in many games (e.g.,
adaptive enemy scaling, hint systems) are implicitly Flow-optimization
mechanisms. The challenge-skill ratio has since become a standard
construct for computational Flow modeling.

The present work operationalizes Csikszentmihalyi's model through a
rule-based classifier that maps the challenge\_skill\_ratio (computed
from behavioral telemetry) to one of four states: FLOW (ratio in
{[}0.75, 1.30{]}), ANXIETY (ratio \textgreater{} 1.30), BOREDOM (ratio
\textless{} 0.75), and APATHY (a joint condition of high idle time, low
dialogue engagement, and low challenge-skill ratio). These thresholds
are informed by prior work on computational Flow detection (Yannakakis
\& Hallam, 2009) and are detailed formally in Section 3.3.4. Critically,
the reward function of the PPO adaptation engine is directly shaped
around these states, making Flow the central organizing principle of the
adaptation objective.

\subsubsection{2.4 Large Language Models for Game
Narrative}\label{large-language-models-for-game-narrative}

\paragraph{2.4.1 PANGeA (AAAI AIIDE 2024)}\label{pangea-aaai-aiide-2024}

Todd et al.'s (2024) Procedural Artificial Narrative using Generative AI
(PANGeA) represents the most directly relevant antecedent work for the
narrative generation component of the current system. PANGeA employs an
LLM to generate narrative content procedurally, grounded in a structured
game knowledge base delivered through the prompt. The key insight from
PANGeA is that LLMs can produce coherent, in-world narrative content
when given sufficient contextual grounding --- but require careful
prompt engineering and content validation to avoid outputs that break
the fictional frame.

PANGeA does not, however, integrate a real-time player model or an
adaptive policy. Its narrative generation is procedural (driven by game
state) rather than personalized (driven by player psychological state).
The current work extends the PANGeA paradigm by adding the player
modeling and RL adaptation layers that PANGeA lacks.

\paragraph{2.4.2 LIGS (CHI 2025)}\label{ligs-chi-2025}

The LLM-Integrated Game Storytelling (LIGS) system, published at CHI
2025, investigated the integration of LLMs into game narrative pipelines
with a focus on maintaining narrative coherence across extended play
sessions. LIGS introduced mechanisms for tracking narrative state across
sessions and using that state to condition generation, anticipating the
episodic memory approach of the current work.

A notable contribution of LIGS was its evaluation methodology: it
employed player experience questionnaires to compare AI-generated
narrative against human-authored content, finding that players could not
reliably distinguish the two when the LLM was appropriately constrained.
This finding provides empirical support for the feasibility of LLM-based
narrative personalization as a subjectively valid approach.

\paragraph{2.4.3 Prompt Grounding and
RAG}\label{prompt-grounding-and-rag}

The challenge of keeping LLM outputs factually and fictionally
consistent with a specific game world has motivated the adoption of
Retrieval-Augmented Generation (RAG; Lewis et al., 2020) in game
narrative contexts. RAG supplements the LLM prompt with retrieved
passages from a curated knowledge base, effectively providing the model
with relevant context at generation time rather than requiring this
knowledge to be baked into model weights through fine-tuning.

For game narrative, RAG offers a particularly attractive tradeoff: the
lore corpus can be maintained and updated independently of the model,
and retrieval ensures that only contextually relevant lore is included
in the prompt, avoiding context window saturation. The current system
employs the sentence-transformers all-MiniLM-L6-v2 model (Wang et al.,
2020) to embed both the lore corpus and generation queries into a
384-dimensional semantic space, with top-k cosine similarity retrieval
(k = 3) providing the relevant passages.

\subsubsection{2.5 Reinforcement Learning for Game
Adaptation}\label{reinforcement-learning-for-game-adaptation}

\paragraph{2.5.1 Dynamic Difficulty
Adjustment}\label{dynamic-difficulty-adjustment}

Dynamic Difficulty Adjustment (DDA) has been an active research area
since Hunicke and Chapman (2004) demonstrated that adaptive difficulty
improved player retention in Super Mario Bros.~Computational approaches
to DDA have employed rule-based systems, Bayesian models, and neural
networks to infer the appropriate challenge level from behavioral
signals and adjust game parameters accordingly. The predominant
formulation is as a feedback control problem: the challenge level is a
controlled variable, player engagement state is observed, and the
controller minimizes the deviation from a target engagement state.

The present work extends the DDA paradigm in two important respects.
First, the adaptation target is not mechanical difficulty (enemy health,
spawn rate, puzzle complexity) but narrative content: the tone, urgency,
and type of narrative delivery. Second, the adaptation mechanism is a
learned policy rather than a hand-crafted controller. These extensions
transform DDA from a mechanical tuning problem into a narrative and
psychological design problem.

\paragraph{2.5.2 PPO and
Stable-Baselines3}\label{ppo-and-stable-baselines3}

Proximal Policy Optimization (PPO; Schulman et al., 2017) is a policy
gradient reinforcement learning algorithm that achieves reliable
performance across a broad range of environments through a clipped
surrogate objective that prevents destructively large policy updates.
PPO has become a standard baseline in both academic RL research and
practical game-playing applications, owing to its relative simplicity,
sample efficiency, and stability.

Stable-Baselines3 (Raffin et al., 2021) provides production-quality
PyTorch implementations of PPO and other RL algorithms, with interfaces
compatible with the Gymnasium (formerly OpenAI Gym) environment
specification. The current work employs Stable-Baselines3 v2.x with the
MlpPolicy architecture, training over 200,000 timesteps in a custom
Gymnasium environment that simulates narrative adaptation episodes. The
use of Stable-Baselines3 ensures reproducibility and provides a
well-tested optimization implementation.

\subsubsection{2.6 Episodic Memory for LLMs (EM-LLM, ICLR
2025)}\label{episodic-memory-for-llms-em-llm-iclr-2025}

Fountas et al.~(2024) introduced EM-LLM, an architecture for endowing
LLMs with human-like episodic memory. The system partitions a long
context stream into discrete memory events, scores them by importance
and surprise, and selectively retrieves the most relevant events for
inclusion in future prompts. This approach enables a form of infinite
context through principled compression: rather than attending to every
prior interaction, the model recalls the events most salient to the
current context.

The current work incorporates a lightweight episodic memory inspired by
EM-LLM. The EpisodicMemory class maintains up to ten narrative events,
each scored by importance, and evicts the least important event when
capacity is exceeded. At generation time, the top five events by
importance score are retrieved and formatted as a context block in the
LLM prompt. This ensures narrative continuity --- the generated content
acknowledges prior story events --- without saturating the context
window.

\subsubsection{2.7 Research Gap and
Positioning}\label{research-gap-and-positioning}

Table 1 summarizes the key related works and their coverage of the
dimensions relevant to this study.

\textbf{Table 1 --- Summary of related work}

\begin{longtable}[]{@{}
  >{\raggedright\arraybackslash}p{(\columnwidth - 12\tabcolsep) * \real{0.1429}}
  >{\raggedright\arraybackslash}p{(\columnwidth - 12\tabcolsep) * \real{0.1429}}
  >{\raggedright\arraybackslash}p{(\columnwidth - 12\tabcolsep) * \real{0.1429}}
  >{\raggedright\arraybackslash}p{(\columnwidth - 12\tabcolsep) * \real{0.1429}}
  >{\raggedright\arraybackslash}p{(\columnwidth - 12\tabcolsep) * \real{0.1429}}
  >{\raggedright\arraybackslash}p{(\columnwidth - 12\tabcolsep) * \real{0.1429}}
  >{\raggedright\arraybackslash}p{(\columnwidth - 12\tabcolsep) * \real{0.1429}}@{}}
\toprule\noalign{}
\begin{minipage}[b]{\linewidth}\raggedright
Work
\end{minipage} & \begin{minipage}[b]{\linewidth}\raggedright
Player Modeling
\end{minipage} & \begin{minipage}[b]{\linewidth}\raggedright
Flow-Based Reward
\end{minipage} & \begin{minipage}[b]{\linewidth}\raggedright
LLM Generation
\end{minipage} & \begin{minipage}[b]{\linewidth}\raggedright
RAG
\end{minipage} & \begin{minipage}[b]{\linewidth}\raggedright
Episodic Memory
\end{minipage} & \begin{minipage}[b]{\linewidth}\raggedright
Complete Loop
\end{minipage} \\
\midrule\noalign{}
\endhead
\bottomrule\noalign{}
\endlastfoot
PANGeA (Todd et al., 2024) & No & No & Yes & Partial & No & No \\
LIGS (CHI 2025) & No & No & Yes & No & Partial & No \\
DDA Survey (Hunicke \& Chapman, 2004) & Behavioral & Yes (implicit) & No
& No & No & No \\
player2vec (2024) & Deep embedding & No & No & No & No & No \\
EM-LLM (Fountas et al., 2024) & No & No & Yes & No & Yes & No \\
\textbf{This work} & \textbf{Hexad + Flow} & \textbf{Yes (explicit)} &
\textbf{Yes (RAG-grounded)} & \textbf{Yes} & \textbf{Yes} &
\textbf{Yes} \\
\end{longtable}

No prior work has integrated continuous Hexad profiling derived from
telemetry, Flow-based RL reward, RAG-grounded LLM generation, and
episodic session memory into a single closed-loop system. The present
work addresses this gap.

\begin{center}\rule{0.5\linewidth}{0.5pt}\end{center}

\subsection{Chapter 3: Research
Methodology}\label{chapter-3-research-methodology}

\subsubsection{3.1 Research Design
Overview}\label{research-design-overview}

The research adopts a mixed-methods design combining a controlled
between-subjects experiment with a qualitative interview study. The
experimental component enables causal inference about the effect of the
AI-driven personalization system on player experience, while the
qualitative component captures the lived experience of personalization
that quantitative measures may not fully represent.

The between-subjects design was chosen over a within-subjects design to
eliminate order effects and carry-over: a player who first experiences
adaptive narrative and then static narrative (or vice versa) would bring
expectations and memory from the first condition into the second. With a
between-subjects design, each participant experiences only one
condition, removing this contamination at the cost of requiring a larger
sample and careful matching.

The experimental condition (Experimental) deploys the full five-module
AI system. The control condition (Control) replaces the Narrative
Generation Service and Adaptation Engine with a StaticNarrativeContent
ScriptableObject that delivers pre-authored narrative content in
response to the same game triggers, with no player-adaptive behavior.
All other aspects of the game session --- environment, objectives, NPC
placement, session duration --- are identical across conditions.

\subsubsection{3.2 System Architecture}\label{system-architecture}

The system is organized as five modules in a closed-loop pipeline.
Figure 1 below describes the architecture and the data flows between
components.

\begin{Shaded}
\begin{Highlighting}[]
\NormalTok{flowchart LR}
\NormalTok{    A["Unity Client\textbackslash{}n(PlayerDataLogger)"] {-}{-}\textgreater{}|WebSocket\textbackslash{}ntelemetry batch\textbackslash{}nevery 5s| B["FastAPI Backend\textbackslash{}n/telemetry/batch"]}
\NormalTok{    B {-}{-}\textgreater{} C["Player Modeling\textbackslash{}nService\textbackslash{}n(FeatureExtractor +\textbackslash{}nHexadProfiler +\textbackslash{}nFlowClassifier)"]}
\NormalTok{    C {-}{-}\textgreater{} D["Adaptation Engine\textbackslash{}n(NarrativeAgent / PPO)"]}
\NormalTok{    D {-}{-}\textgreater{}|NarrativeAction| E["Narrative Generation\textbackslash{}nService\textbackslash{}n(LoreRetriever +\textbackslash{}nEpisodicMemory +\textbackslash{}nPromptBuilder +\textbackslash{}nOllamaClient)"]}
\NormalTok{    E {-}{-}\textgreater{}|NarrativeContent JSON| F["Unity Client\textbackslash{}n(NarrativeManager +\textbackslash{}nContentInjector)"]}
\NormalTok{    F {-}{-}\textgreater{}|in{-}game dialogue\textbackslash{}ndescription\textbackslash{}nquest update| G["Player"]}
\NormalTok{    G {-}{-}\textgreater{}|play behavior| A}
\end{Highlighting}
\end{Shaded}

\textbf{Figure 1 --- High-level closed-loop architecture}

The pipeline operates as follows. The Unity PlayerDataLogger aggregates
behavioral signals over five-second windows and transmits a
TelemetryBatch to the FastAPI backend via HTTP POST (with WebSocket
support for lower-latency streaming). The Player Modeling Service
processes incoming batches: the FeatureExtractor computes eleven
normalized features, the HexadProfiler derives the six-dimensional
continuous Hexad profile, and the FlowClassifier assigns a Flow state.
The resulting PlayerModel is passed to the NarrativeAgent, which selects
a NarrativeAction using either the cold-start heuristic (first 120
seconds) or the trained PPO policy. The Narrative Generation Service
uses the selected action to query the LoreRetriever (RAG), consults the
EpisodicMemory, assembles the full prompt through the PromptBuilder, and
requests generation from the Ollama-hosted LLM. The resulting
NarrativeContent is returned to the Unity client and injected at the
appropriate in-game display point.

\subsubsection{3.3 Player Modeling Module}\label{player-modeling-module}

\paragraph{3.3.1 Telemetry Collection}\label{telemetry-collection}

The telemetry system captures 23 behavioral fields per five-second
window. Table 2 lists the complete schema, as defined in
\texttt{backend/models/telemetry.py}.

\textbf{Table 2 --- TelemetryBatch schema (23 fields)}

\begin{longtable}[]{@{}
  >{\raggedright\arraybackslash}p{(\columnwidth - 4\tabcolsep) * \real{0.3333}}
  >{\raggedright\arraybackslash}p{(\columnwidth - 4\tabcolsep) * \real{0.3333}}
  >{\raggedright\arraybackslash}p{(\columnwidth - 4\tabcolsep) * \real{0.3333}}@{}}
\toprule\noalign{}
\begin{minipage}[b]{\linewidth}\raggedright
Field
\end{minipage} & \begin{minipage}[b]{\linewidth}\raggedright
Type
\end{minipage} & \begin{minipage}[b]{\linewidth}\raggedright
Description
\end{minipage} \\
\midrule\noalign{}
\endhead
\bottomrule\noalign{}
\endlastfoot
player\_id & str & Unique player identifier \\
session\_id & str & Unique session identifier \\
window\_start & float & Window start time (Unix timestamp) \\
window\_end & float & Window end time (Unix timestamp) \\
kills & int & Enemy kills in window \\
deaths & int & Player deaths in window \\
damage\_taken & float & Total damage received \\
damage\_dealt & float & Total damage inflicted \\
abilities\_used & int & Ability activations \\
abilities\_hit & int & Abilities that connected \\
objectives\_completed & int & Objectives completed \\
objectives\_attempted & int & Objectives attempted \\
tiles\_explored & int & Map tiles visited \\
total\_tiles & int & Total traversable tiles (default 100) \\
npc\_interactions\_voluntary & int & Voluntary (unsolicited) NPC
interactions \\
dialogue\_lines\_shown & int & Dialogue lines displayed \\
dialogue\_lines\_skipped & int & Dialogue lines skipped by player \\
lore\_items\_read & int & Lore collectibles read \\
lore\_items\_found & int & Lore collectibles discovered \\
idle\_seconds & float & Seconds with no player input \\
session\_elapsed & float & Total elapsed session time (seconds) \\
current\_location & str & Scene/zone identifier \\
current\_objective & str & Active quest objective identifier \\
\end{longtable}

The five-second window granularity balances responsiveness (enabling
sub-minute adaptation cycles) against noise sensitivity (single-frame
events are smoothed over the window). In practice, the system
accumulates a sliding window of recent batches before computing the
feature vector, allowing multi-window averages to further smooth
transient spikes.

\paragraph{3.3.2 Feature Extraction (11
Features)}\label{feature-extraction-11-features}

The FeatureExtractor aggregates one or more TelemetryBatch records into
eleven normalized scalar features, each clipped to {[}0, 1{]} (or {[}0,
3{]} for the challenge\_skill\_ratio). Table 3 presents the computation
formulae.

\textbf{Table 3 --- Eleven extracted features and computation formulae}

\begin{longtable}[]{@{}
  >{\raggedright\arraybackslash}p{(\columnwidth - 4\tabcolsep) * \real{0.3333}}
  >{\raggedright\arraybackslash}p{(\columnwidth - 4\tabcolsep) * \real{0.3333}}
  >{\raggedright\arraybackslash}p{(\columnwidth - 4\tabcolsep) * \real{0.3333}}@{}}
\toprule\noalign{}
\begin{minipage}[b]{\linewidth}\raggedright
Feature
\end{minipage} & \begin{minipage}[b]{\linewidth}\raggedright
Formula
\end{minipage} & \begin{minipage}[b]{\linewidth}\raggedright
Notes
\end{minipage} \\
\midrule\noalign{}
\endhead
\bottomrule\noalign{}
\endlastfoot
kd\_ratio & clip(kills / max(deaths, 1) / 5.0, 0, 1) & Normalized at a
5:1 K/D ceiling \\
objective\_rate & clip(obj\_completed / max(obj\_attempted, 1), 0, 1) &
Proportion of attempted objectives completed \\
explore\_rate & clip(tiles\_explored / total\_tiles, 0, 1) & Map
coverage fraction \\
dialogue\_engage\_rate & clip(1 - d\_skipped / max(d\_shown, 1), 0, 1) &
1 = reads all dialogue \\
ability\_accuracy & clip(ab\_hit / max(ab\_used, 1), 0, 1) & Combat
precision \\
damage\_taken\_norm & clip(damage\_taken / (elapsed\_minutes + 1) / 200,
0, 1) & DPM normalized \\
idle\_ratio & clip(idle\_seconds / elapsed, 0, 1) & Fraction of session
idle \\
lore\_engage\_rate & clip(lore\_read / max(lore\_found, 1), 0, 1) & Lore
engagement depth \\
skill\_score & 0.35\emph{kd + 0.30}obj + 0.20\emph{acc + 0.15}(1-idle) &
Composite skill index \\
challenge\_score & 0.50\emph{dmg\_norm + 0.30}(1-obj\_rate) +
0.20*(deaths/(deaths+kills+1)) & Composite challenge index \\
challenge\_skill\_ratio & clip(challenge\_score / max(skill\_score,
1e-6), 0, 3) & Core Flow signal \\
\end{longtable}

The skill\_score composite weights combat effectiveness (kd\_ratio,
ability\_accuracy) most heavily, followed by objective completion and
attention (inverse idle\_ratio). The challenge\_score composite weights
damage pressure most heavily, followed by objective failure rate and
death rate. Dividing challenge by skill yields the
challenge\_skill\_ratio, which operationalizes Csikszentmihalyi's (1990)
challenge-skill balance axis.

\paragraph{3.3.3 Hexad Profile Computation (6
Dimensions)}\label{hexad-profile-computation-6-dimensions}

The HexadProfiler computes six continuous scores in {[}0, 1{]} from
accumulated telemetry, each representing the strength of a Hexad player
type signal. Table 4 shows the computation weights.

\textbf{Table 4 --- Hexad dimension computation weights}

\begin{longtable}[]{@{}
  >{\raggedright\arraybackslash}p{(\columnwidth - 4\tabcolsep) * \real{0.3333}}
  >{\raggedright\arraybackslash}p{(\columnwidth - 4\tabcolsep) * \real{0.3333}}
  >{\raggedright\arraybackslash}p{(\columnwidth - 4\tabcolsep) * \real{0.3333}}@{}}
\toprule\noalign{}
\begin{minipage}[b]{\linewidth}\raggedright
Dimension
\end{minipage} & \begin{minipage}[b]{\linewidth}\raggedright
Formula
\end{minipage} & \begin{minipage}[b]{\linewidth}\raggedright
Behavioral Interpretation
\end{minipage} \\
\midrule\noalign{}
\endhead
\bottomrule\noalign{}
\endlastfoot
Achiever & 0.70 * obj\_rate + 0.30 * norm(kills, 20) & High objective
completion + combat activity \\
Explorer & 0.50 * explore\_rate + 0.50 * lore\_rate & Map coverage +
lore engagement \\
Socializer & 0.60 * norm(voluntary\_npc, 10) + 0.40 * dialogue\_engage &
Seeks NPCs + reads dialogue \\
Free Spirit & 0.50 * explore\_rate + 0.50 * (1 - obj\_rate) & Explores
without following objectives \\
Disruptor & 0.65 * norm(kills, 20) + 0.15 * norm(deaths, 10) + 0.20 *
norm(damage\_dealt, 1000) & Combat-heavy, aggressive \\
Philanthropist & 0.60 * dialogue\_engage + 0.40 * lore\_rate & Engages
with narrative and lore \\
\end{longtable}

The mapping from behavioral signals to Hexad dimensions was established
through a principled alignment process guided by the original Hexad
motivation descriptions (Tondello et al., 2016): Achievers seek mastery
and progress (captured by objective completion and combat
effectiveness); Explorers seek discovery (captured by map and lore
coverage); Socializers seek connection (captured by voluntary NPC
interaction and dialogue engagement); Free Spirits resist imposed
structure (captured by exploration uncorrelated with objective
completion); Disruptors challenge systems (captured by aggressive combat
behavior); Philanthropists give and share (captured by engagement with
narrative and lore content that would benefit future players'
understanding).

All profiles are initialized to 0.5 (neutral) in the absence of data,
preventing cold-start artifacts from biasing early narrative selection.

\paragraph{3.3.4 Flow State Classification (MDP
States)}\label{flow-state-classification-mdp-states}

The FlowClassifier implements a decision-tree classifier over the
challenge\_skill\_ratio, idle\_ratio, and dialogue\_engage\_rate
features. Table 5 formalizes the rules.

\textbf{Table 5 --- Flow state classification rules (evaluated in
priority order)}

\begin{longtable}[]{@{}
  >{\raggedright\arraybackslash}p{(\columnwidth - 6\tabcolsep) * \real{0.2500}}
  >{\raggedright\arraybackslash}p{(\columnwidth - 6\tabcolsep) * \real{0.2500}}
  >{\raggedright\arraybackslash}p{(\columnwidth - 6\tabcolsep) * \real{0.2500}}
  >{\raggedright\arraybackslash}p{(\columnwidth - 6\tabcolsep) * \real{0.2500}}@{}}
\toprule\noalign{}
\begin{minipage}[b]{\linewidth}\raggedright
Priority
\end{minipage} & \begin{minipage}[b]{\linewidth}\raggedright
Condition
\end{minipage} & \begin{minipage}[b]{\linewidth}\raggedright
State
\end{minipage} & \begin{minipage}[b]{\linewidth}\raggedright
Confidence
\end{minipage} \\
\midrule\noalign{}
\endhead
\bottomrule\noalign{}
\endlastfoot
1 (highest) & idle \textgreater{} 0.40 AND dialogue \textless{} 0.30 AND
ratio \textless{} 0.60 & APATHY & 0.70 \\
2 & ratio \textless{} 0.75 & BOREDOM & 0.60 + (0.75 - ratio) * 0.80 \\
3 & ratio \textgreater{} 1.30 & ANXIETY & 0.60 + (ratio - 1.30) *
0.40 \\
4 (default) & 0.75 ≤ ratio ≤ 1.30 & FLOW & max(0.55, 1.0 - \\
\end{longtable}

APATHY is classified first because it represents a qualitatively
distinct disengagement pattern (high inactivity + low narrative
interaction + low challenge) that would otherwise be misclassified as
BOREDOM on the ratio criterion alone. The confidence scoring for BOREDOM
and ANXIETY is proportional to the distance from the FLOW band boundary,
reflecting that states further from the boundary are classified with
higher certainty. FLOW confidence is highest when the ratio is closest
to 1.0 (perfect challenge-skill balance) and decreases toward the
boundary thresholds.

The four Flow states directly correspond to the MDP state space of the
adaptation engine (Section 3.5), establishing the tight coupling between
the player modeling and RL components.

\subsubsection{3.4 Narrative Generation
Module}\label{narrative-generation-module}

The narrative generation module receives a NarrativeAction from the
adaptation engine and produces a NarrativeContent object --- a
structured JSON payload containing the type of content (dialogue,
description, or quest update), the generated text, the speaker (if
dialogue), the emotional tone, and any lore references. Figure 2
illustrates the internal data flow.

\begin{Shaded}
\begin{Highlighting}[]
\NormalTok{flowchart TD}
\NormalTok{    A[NarrativeAction\textbackslash{}nfrom PPO agent] {-}{-}\textgreater{} B[LoreRetriever\textbackslash{}nRAG query]}
\NormalTok{    C[PlayerModel] {-}{-}\textgreater{} B}
\NormalTok{    C {-}{-}\textgreater{} D[PromptBuilder]}
\NormalTok{    B {-}{-}\textgreater{}|top{-}3 lore chunks| D}
\NormalTok{    E[EpisodicMemory] {-}{-}\textgreater{}|top{-}5 events| D}
\NormalTok{    D {-}{-}\textgreater{}|system prompt\textbackslash{}n+ user prompt| F[OllamaClient\textbackslash{}nPhi{-}3.5 Mini]}
\NormalTok{    F {-}{-}\textgreater{}|JSON response| G\{Parse \&\textbackslash{}nValidate\}}
\NormalTok{    G {-}{-}\textgreater{}|valid| H[NarrativeContent]}
\NormalTok{    G {-}{-}\textgreater{}|fallback| I[Static fallback\textbackslash{}ntext]}
\end{Highlighting}
\end{Shaded}

\textbf{Figure 2 --- Narrative generation pipeline}

\paragraph{3.4.1 RAG Pipeline}\label{rag-pipeline}

The LoreRetriever implements in-memory semantic search using the
sentence-transformers all-MiniLM-L6-v2 model, which produces
384-dimensional L2-normalized embeddings. The lore corpus is pre-indexed
at startup by chunking the world, character, and quest lore documents by
paragraph (double-newline delimiters) and embedding each chunk. At
retrieval time, the query is embedded and cosine similarities are
computed against all stored embeddings through a dot product (since all
vectors are unit-normalized, cosine similarity equals the dot product).
The top k = 3 chunks by similarity are returned and formatted as bullet
points in the LLM user prompt.

This design was chosen over a persistent vector database (e.g.,
ChromaDB, Pinecone) for two reasons: Python 3.14 compatibility issues
with ChromaDB's Pydantic v1 dependency, and the relatively small lore
corpus size (three files, approximately 30 paragraphs total) which
renders the overhead of a full vector database unnecessary. In-memory
retrieval at this corpus size is effectively instantaneous and imposes
no latency penalty on the adaptation cycle.

\paragraph{3.4.2 Episodic Session Memory}\label{episodic-session-memory}

The EpisodicMemory class maintains a capped list of MemoryEvent records,
each comprising a text description, an importance score in {[}0, 1{]},
and a timestamp. When the list exceeds the maximum capacity of ten
events, it is sorted by importance in descending order and truncated to
the top ten. This importance-weighted eviction ensures that the most
narratively significant events (e.g., the player's first encounter with
a major NPC, a key plot revelation) are retained even as minor events
accumulate.

At prompt construction time, the top five events by importance score are
selected and formatted as a context block:

\begin{verbatim}
KEY NARRATIVE EVENTS THIS SESSION:
- Player spoke with Captain Aldric at the village gate
- Player descended into Dungeon Room A
- Player discovered the Warden's Log fragment
\end{verbatim}

This mechanism, inspired by EM-LLM (Fountas et al., 2024), provides the
LLM with a compressed narrative history that enables references to past
events, preventing the temporal inconsistency that would arise from
generating each narrative beat in isolation.

\paragraph{3.4.3 Dynamic Prompt
Construction}\label{dynamic-prompt-construction}

The PromptBuilder assembles a two-message prompt (system + user) for the
Ollama API. The system prompt encodes the LLM's role (narrative writer
for a dark fantasy RPG called Eryndal), output length constraints (2-4
sentences for dialogue, 1-2 for descriptions), format requirements
(valid JSON only, specific schema), and the emotional tone mapped from
the current Flow state:

\begin{itemize}
\tightlist
\item
  ANXIETY: ``urgent and supportive --- the player is struggling''
\item
  BOREDOM: ``exciting and surprising --- inject energy''
\item
  FLOW: ``immersive and atmospheric --- maintain the experience''
\item
  APATHY: ``intriguing and mysterious --- re-engage the player''
\end{itemize}

The user prompt injects: the player's current Flow state and
challenge-skill ratio; the current activity (IN\_COMBAT, EXPLORING,
IN\_DIALOGUE, IDLE); session elapsed time; the Hexad profile (Explorer,
Socializer, Achiever, Disruptor percentages); the selected
NarrativeAction with a brief definition of each of the eight actions;
the current location and active quest stage; the three retrieved lore
chunks; the episodic memory context; and explicit content constraints
(no real-world references, no anachronisms, no fourth-wall breaks,
character voice descriptions for each of the three NPCs). This rich
contextual grounding aims to constrain the generation to in-world,
narratively appropriate content.

\paragraph{3.4.4 LLM Integration (Ollama / Phi-3.5
Mini)}\label{llm-integration-ollama-phi-3.5-mini}

The OllamaClient sends the assembled prompt to a locally hosted Ollama
instance using the /api/chat endpoint. The primary model is Phi-3.5 Mini
(3.8B parameters; Microsoft, 2024), with Llama 3.2 3B as a configured
fallback. The generation parameters are: temperature 0.75, top\_p 0.90,
maximum tokens 200, JSON format mode enabled. The 8-second HTTP timeout
ensures that a slow or stalled generation does not block the main game
loop.

When the LLM response is malformed JSON, contains a \{``fallback'':
true\} signal, or the request times out, the OllamaClient returns a
pre-authored fallback string corresponding to the requested
NarrativeAction. These fallback strings are minimal, tonally appropriate
responses (e.g., for PROVIDE\_GUIDANCE: ``The path ahead is treacherous.
Tread carefully.'') that maintain narrative coherence without exposing
the player to system failure artifacts.

\subsubsection{3.5 Adaptation Engine (MDP +
PPO)}\label{adaptation-engine-mdp-ppo}

\paragraph{3.5.1 MDP Formulation}\label{mdp-formulation}

The narrative adaptation problem is formalized as a Markov Decision
Process (S, A, T, R, γ) implemented as a Gymnasium environment.

\textbf{State space S:} A Box observation space of shape (7,) with all
values in {[}0, 1{]}. The seven dimensions encode:

\begin{longtable}[]{@{}
  >{\raggedright\arraybackslash}p{(\columnwidth - 4\tabcolsep) * \real{0.3333}}
  >{\raggedright\arraybackslash}p{(\columnwidth - 4\tabcolsep) * \real{0.3333}}
  >{\raggedright\arraybackslash}p{(\columnwidth - 4\tabcolsep) * \real{0.3333}}@{}}
\toprule\noalign{}
\begin{minipage}[b]{\linewidth}\raggedright
Index
\end{minipage} & \begin{minipage}[b]{\linewidth}\raggedright
Signal
\end{minipage} & \begin{minipage}[b]{\linewidth}\raggedright
Description
\end{minipage} \\
\midrule\noalign{}
\endhead
\bottomrule\noalign{}
\endlastfoot
0 & flow\_score & Encoded Flow state (FLOW=1.0, BOREDOM=0.3,
ANXIETY=0.2, APATHY=0.1) \\
1 & challenge\_skill\_ratio\_norm & challenge\_skill\_ratio / 3.0 \\
2 & hexad\_explorer & Explorer dimension score \\
3 & hexad\_socializer & Socializer dimension score \\
4 & hexad\_achiever & Achiever dimension score \\
5 & hexad\_disruptor & Disruptor dimension score \\
6 & session\_progress & session\_elapsed / 1800.0 (normalized to 30-min
max) \\
\end{longtable}

\textbf{Action space A:} A Discrete(8) action space corresponding to the
eight NarrativeAction values. Table 6 summarizes the actions and their
intended narrative effects.

\textbf{Table 6 --- Eight narrative actions with intended effects}

\begin{longtable}[]{@{}
  >{\raggedright\arraybackslash}p{(\columnwidth - 4\tabcolsep) * \real{0.3333}}
  >{\raggedright\arraybackslash}p{(\columnwidth - 4\tabcolsep) * \real{0.3333}}
  >{\raggedright\arraybackslash}p{(\columnwidth - 4\tabcolsep) * \real{0.3333}}@{}}
\toprule\noalign{}
\begin{minipage}[b]{\linewidth}\raggedright
Index
\end{minipage} & \begin{minipage}[b]{\linewidth}\raggedright
Action
\end{minipage} & \begin{minipage}[b]{\linewidth}\raggedright
Intended Effect
\end{minipage} \\
\midrule\noalign{}
\endhead
\bottomrule\noalign{}
\endlastfoot
0 & LOWER\_STAKES & Reduce narrative tension; provide respite \\
1 & RAISE\_STAKES & Increase narrative urgency and weight \\
2 & ADD\_MYSTERY & Introduce a curiosity hook or unexplained element \\
3 & ADD\_HUMOR & Provide comic relief to reduce tension \\
4 & PROVIDE\_GUIDANCE & Give the struggling player directional narrative
support \\
5 & INCREASE\_URGENCY & Re-energize a bored or disengaged player \\
6 & LORE\_REWARD & Deliver a lore discovery payoff for curious
players \\
7 & NO\_CHANGE & Suppress narrative injection; maintain current
experience \\
\end{longtable}

\textbf{Transition dynamics T:} The simulated training environment
models stochastic Flow state transitions. When the agent selects a
``helpful'' action for the current state (ANXIETY→\{PROVIDE\_GUIDANCE,
LOWER\_STAKES\}, BOREDOM→\{INCREASE\_URGENCY, RAISE\_STAKES,
ADD\_MYSTERY\}, APATHY→\{ADD\_MYSTERY, LORE\_REWARD\},
FLOW→\{NO\_CHANGE\}), the environment transitions to FLOW with
probability 0.75 or remains in the current state with probability 0.25.
This stochasticity prevents the policy from memorizing a deterministic
mapping and encourages robust adaptation.

\textbf{Discount factor γ:} 0.95, reflecting moderate preference for
near-term reward while maintaining long-horizon planning.

Figure 4 illustrates the MDP state transitions.

\begin{Shaded}
\begin{Highlighting}[]
\NormalTok{stateDiagram{-}v2}
\NormalTok{    FLOW {-}{-}\textgreater{} BOREDOM : helpful action fails (p=0.10)}
\NormalTok{    FLOW {-}{-}\textgreater{} FLOW : NO\_CHANGE or helpful succeeds (p=0.90)}
\NormalTok{    BOREDOM {-}{-}\textgreater{} FLOW : helpful action succeeds (p=0.75)}
\NormalTok{    BOREDOM {-}{-}\textgreater{} BOREDOM : unhelpful or fails (p=0.25)}
\NormalTok{    ANXIETY {-}{-}\textgreater{} FLOW : PROVIDE\_GUIDANCE / LOWER\_STAKES (p=0.75)}
\NormalTok{    ANXIETY {-}{-}\textgreater{} ANXIETY : unhelpful or fails (p=0.25)}
\NormalTok{    APATHY {-}{-}\textgreater{} FLOW : ADD\_MYSTERY / LORE\_REWARD (p=0.75)}
\NormalTok{    APATHY {-}{-}\textgreater{} APATHY : unhelpful or fails (p=0.25)}
\end{Highlighting}
\end{Shaded}

\textbf{Figure 4 --- MDP state transition diagram}

\paragraph{3.5.2 Reward Function Design}\label{reward-function-design}

The reward function is defined over (prev\_state, new\_state, action,
steps\_in\_state, last\_3\_actions) and comprises five components:

\textbf{Table 8 --- Reward function components}

\begin{longtable}[]{@{}
  >{\raggedright\arraybackslash}p{(\columnwidth - 4\tabcolsep) * \real{0.3333}}
  >{\raggedright\arraybackslash}p{(\columnwidth - 4\tabcolsep) * \real{0.3333}}
  >{\raggedright\arraybackslash}p{(\columnwidth - 4\tabcolsep) * \real{0.3333}}@{}}
\toprule\noalign{}
\begin{minipage}[b]{\linewidth}\raggedright
Component
\end{minipage} & \begin{minipage}[b]{\linewidth}\raggedright
Value
\end{minipage} & \begin{minipage}[b]{\linewidth}\raggedright
Trigger
\end{minipage} \\
\midrule\noalign{}
\endhead
\bottomrule\noalign{}
\endlastfoot
State reward (FLOW) & +1.0 & new\_state == FLOW \\
State reward (BOREDOM) & -0.3 & new\_state == BOREDOM \\
State reward (ANXIETY) & -0.5 & new\_state == ANXIETY \\
State reward (APATHY) & -0.8 & new\_state == APATHY \\
Step penalty & -0.05 & Every step \\
Flow streak bonus & +0.20 & new\_state == FLOW AND steps\_in\_state
\textgreater{} 2 \\
Transition bonus & +0.30 & prev in \{ANXIETY, BOREDOM\} AND new\_state
== FLOW \\
Repetition penalty & -0.15 & All 3 recent actions are identical \\
\end{longtable}

The state rewards reflect the severity ordering of non-Flow states:
APATHY is penalized most heavily because it represents the most
concerning disengagement (passive inactivity rather than active
frustration or excitement). The step penalty provides a constant
pressure toward action efficiency. The Flow streak bonus reinforces
sustained optimal experience. The transition bonus provides additional
shaping reward for the key recovery transitions. The repetition penalty
discourages the degenerate policy of repeatedly applying the same action
regardless of state.

The total reward at any step is bounded approximately in {[}-1.10,
+1.45{]}, providing sufficient gradient signal for stable PPO learning.

\paragraph{3.5.3 Cold-Start Heuristic}\label{cold-start-heuristic}

For the first 120 seconds of a session, or when the PPO model is
unavailable, the NarrativeAgent falls back to a deterministic heuristic
policy. Figure 5 illustrates this.

\begin{Shaded}
\begin{Highlighting}[]
\NormalTok{flowchart TD}
\NormalTok{    A\{session\_elapsed\textbackslash{}n\textless{} 120s?\} {-}{-}\textgreater{}|Yes| B\{Flow State\}}
\NormalTok{    B {-}{-}\textgreater{}|ANXIETY| C[PROVIDE\_GUIDANCE]}
\NormalTok{    B {-}{-}\textgreater{}|BOREDOM| D[INCREASE\_URGENCY]}
\NormalTok{    B {-}{-}\textgreater{}|APATHY| E[ADD\_MYSTERY]}
\NormalTok{    B {-}{-}\textgreater{}|FLOW| F[NO\_CHANGE]}
\NormalTok{    A {-}{-}\textgreater{}|No| G[PPO Policy]}
\end{Highlighting}
\end{Shaded}

\textbf{Figure 5 --- Cold-start heuristic decision tree}

The 120-second threshold was chosen to allow the player sufficient time
to engage with the game world before the RL policy takes effect, giving
the player modeling module time to accumulate meaningful telemetry.
During this window, the heuristic maps directly from Flow state to the
most intuitively appropriate narrative action: guidance for struggling
players, urgency for bored players, mystery for disengaged players, and
no change for players already in Flow.

\subsubsection{3.6 Unity Plugin Design}\label{unity-plugin-design}

The Unity plugin comprises three C\# components: PlayerDataLogger,
NarrativeManager, and ContentInjector.

\paragraph{3.6.1 PlayerDataLogger}\label{playerdatalogger}

PlayerDataLogger is a MonoBehaviour that maintains running accumulators
for all 23 telemetry fields, firing on Update() at 60fps. Every five
seconds, it serializes the accumulated values into a TelemetryBatch JSON
object and dispatches it to the FastAPI backend via an async HTTP POST
to the /telemetry/batch endpoint. It also exposes a public API for game
systems to register events: \texttt{RecordKill()},
\texttt{RecordDeath()}, \texttt{RecordDialogueShown()},
\texttt{RecordDialogueSkipped()}, \texttt{RecordLoreFound()},
\texttt{RecordLoreRead()}, \texttt{RecordNPCInteraction()},
\texttt{RecordObjectiveAttempted()},
\texttt{RecordObjectiveCompleted()}. This event-based registration model
decouples the telemetry system from specific game mechanics, making the
plugin applicable to any Unity game that exposes these event hooks.

\paragraph{3.6.2 NarrativeManager and
ContentInjector}\label{narrativemanager-and-contentinjector}

NarrativeManager is responsible for requesting narrative content from
the backend. It polls the /narrative/generate endpoint following each
telemetry transmission, receiving a NarrativeContent JSON object that it
deserializes and routes to the ContentInjector.

ContentInjector is a scene-aware component that maps
NarrativeContentType to the appropriate UI mechanism: dialogue content
is routed to the DialoguePanel (a world-space UI overlay above the
relevant NPC); description content is displayed in the DescriptionPanel
(a contextual text panel); quest updates are pushed to the QuestLog and
the QuestUpdateBanner. The injector implements a content queue with a
3-second minimum spacing between consecutive injections to prevent
narrative overcrowding during high-activity moments.

For the Control condition, the NarrativeManager is replaced by a
StaticNarrativeManager that reads from a StaticNarrativeContent
ScriptableObject and fires narrative events at scripted triggers (e.g.,
entering a new scene, completing an objective) with no player-state
dependency.

\subsubsection{3.7 Testbed Game Design}\label{testbed-game-design}

The testbed is a 2D top-down RPG built in Unity 6 LTS, set in the dark
fantasy world of Eryndal. The world is described in
\texttt{E:\textbackslash{}Projects\textbackslash{}fyp\textbackslash{}v2\textbackslash{}lore\textbackslash{}world.md},
character backstories in
\texttt{E:\textbackslash{}Projects\textbackslash{}fyp\textbackslash{}v2\textbackslash{}lore\textbackslash{}characters.md},
and quest structure in
\texttt{E:\textbackslash{}Projects\textbackslash{}fyp\textbackslash{}v2\textbackslash{}lore\textbackslash{}quests.md}.

The kingdom of Eryndal has fallen into shadow following the Sundering
War of thirty years prior. Ancient seals containing Void-touched
creatures have begun to fail. The game world encompasses three navigable
scenes:

\begin{enumerate}
\def\labelenumi{\arabic{enumi}.}
\item
  \textbf{Ashfen Village Hub} --- the starting area, populated by three
  NPCs: Captain Aldric (a guilt-haunted veteran seeking redemption),
  Mira the Merchant (a dry, pragmatic trader with covert dungeon
  knowledge), and The Stranger (a cryptic, unsettling figure whose
  motives are unknown). The village hub serves as a social and narrative
  hub; players who engage extensively with NPCs here register high
  Socializer and Philanthropist signals.
\item
  \textbf{Dungeon Room A (Entry Hall)} --- a combat-focused area
  patrolled by corrupted guards. Players who clear this room efficiently
  register high Achiever and Disruptor signals; players who explore its
  periphery register Explorer signals.
\item
  \textbf{Dungeon Room B (Mid-Level)} --- a hybrid combat-exploration
  area containing the Warden's Log and environmental lore. The presence
  of lore collectibles creates a natural Explorer/Philanthropist signal
  divergence from the Disruptor signal.
\end{enumerate}

The main quest (The Broken Seal) guides players through all three scenes
across five stages: entering Ashfen, investigating the dungeon entrance,
descending to the mid-level, reaching the seal chamber, and recovering
the Sundering Shard. Two side quests (Mira's Map and The Stranger's
Warning) provide additional narrative depth for players who engage with
the social and mystery elements. A standard play session is
approximately 25-30 minutes for players who complete the main quest and
at least one side quest, providing sufficient telemetry for meaningful
player modeling.

\subsubsection{3.8 Evaluation Design}\label{evaluation-design}

\paragraph{3.8.1 Research Hypotheses}\label{research-hypotheses}

The following directional hypotheses are tested:

\begin{itemize}
\tightlist
\item
  \textbf{H1:} Participants in the Experimental condition will score
  significantly higher on the GEQ Flow subscale than participants in the
  Control condition.
\item
  \textbf{H2:} Participants in the Experimental condition will score
  significantly higher on the GEQ Immersion subscale than participants
  in the Control condition.
\item
  \textbf{H3:} Participants in the Experimental condition will score
  significantly higher on the NPS-3 Narrative Personalization Scale than
  participants in the Control condition.
\end{itemize}

H1 is designated the primary hypothesis because Flow is the direct
optimization target of the adaptation engine, making it the most
theoretically motivated outcome. H2 and H3 are secondary hypotheses
testing related but distinct aspects of the experience.

\paragraph{3.8.2 Participants and
Recruitment}\label{participants-and-recruitment}

A total of N = 50 participants are recruited (25 Experimental, 25
Control) via convenience sampling through university mailing lists and
gaming society channels. Inclusion criteria: aged 18 or over;
self-reported experience playing RPG or action-adventure games; no prior
exposure to the testbed game. Exclusion criteria: significant visual
impairment not corrected to normal; participation in a previous pilot
session.

An a priori power analysis for Welch's independent-samples t-test
targeting d = 0.5 (medium effect), α = 0.05 (two-tailed), and 1 - β =
0.80 yields a required sample size of approximately 51 total
participants (Faul et al., 2007), indicating that N = 50 is marginally
below the strict threshold but remains reasonable given the
between-subjects design, the fact that effect sizes for personalization
interventions in games are often reported at d \textgreater{} 0.5
(Tondello et al., 2016), and the practical constraints of master's-level
research.

Participants are randomly assigned to conditions using a balanced random
allocation procedure to ensure approximately equal group sizes. No block
randomization by demographic is applied, given the expected homogeneity
of the university convenience sample.

\paragraph{3.8.3 Instruments}\label{instruments}

\textbf{Game Experience Questionnaire (GEQ) --- Core Module}

The GEQ (IJsselsteijn et al., 2013) is a 33-item instrument measuring
seven dimensions of game experience on a 0--4 scale. The primary
dependent variable is the GEQ \textbf{Flow subscale} (items 5, 13, 25,
28, 31; items 5 and 13 reverse-scored), which captures the absorption,
loss of temporal awareness, and engagement characteristics of the Flow
state. The secondary dependent variable is the GEQ \textbf{Immersion
subscale} (items 3, 12, 18, 19, 27, 30), which captures transportation
into the game world. Subscale scores are computed as means of scored
items. The GEQ Flow subscale has demonstrated acceptable reliability
(Cronbach's α ≈ 0.79) in prior studies with gaming populations.

\textbf{miniPXI --- Player Experience Inventory (Short Form)}

The miniPXI (Vornhagen et al., 2022) is an 11-item instrument on a 1--7
scale covering six player experience dimensions: Mastery (items 1-2),
Curiosity (items 3-4), Immersion (items 5-6), Autonomy (items 7-8),
Meaning (item 8), and Positive Affect (items 9-11). The total score
(mean of all 11 items) is used as a holistic player experience index.
The miniPXI provides a validated, concise complement to the GEQ that
captures dimensions (e.g., autonomy, curiosity) not covered by the GEQ
alone.

\textbf{Narrative Personalization Scale (NPS-3)}

The NPS-3 is a custom three-item scale developed for this study on a
1--7 Likert scale, with items: (1) ``The story felt tailored to how I
was playing,'' (2) ``The narrative responded to my choices and
actions,'' and (3) ``The game seemed to understand what kind of player I
am.'' The NPS-3 directly operationalizes the perceived personalization
construct (H3). Cronbach's alpha will be reported; if α \textless{}
0.70, items will be analyzed individually. The scale was piloted with n
= 3 participants prior to the main study to check item comprehension.

\textbf{Behavioural Metrics (from telemetry)}

Three behavioral metrics derived from the telemetry log serve as
objective secondary outcome measures: (1) voluntary NPC interactions
(count), reflecting unprompted narrative engagement; (2) dialogue read
ratio (1 - skipped/shown), reflecting narrative attention; and (3)
exploration coverage (tiles\_explored / total\_tiles), reflecting
spatial engagement.

\paragraph{3.8.4 Procedure}\label{procedure}

Each participant session follows the procedure illustrated in Figure 8
and is conducted individually in a controlled laboratory setting.

\begin{enumerate}
\def\labelenumi{\arabic{enumi}.}
\tightlist
\item
  \textbf{Arrival and consent} (5 minutes): participant is briefed on
  the study purpose (described as an investigation of game narrative
  experience, without disclosing the adaptation system), provided the
  participant information sheet, and signs informed consent. Audio
  recording consent for the interview is obtained separately.
\item
  \textbf{Tutorial} (5 minutes): participant completes a brief,
  condition-identical tutorial scene outside the experimental
  environment, familiarizing them with movement, combat, and NPC
  interaction mechanics.
\item
  \textbf{Game session} (25-30 minutes): participant plays the testbed
  game in their assigned condition. The researcher is present but
  silent. Session length is time-limited to 30 minutes to ensure
  cross-condition comparability.
\item
  \textbf{Questionnaires} (10 minutes): immediately following the
  session, the participant completes the GEQ Core Module, miniPXI, and
  NPS-3 in paper form.
\item
  \textbf{Semi-structured interview} (20 minutes): the researcher
  conducts the interview following the guide in Appendix D. The session
  is audio recorded with participant consent.
\item
  \textbf{Debrief} (5 minutes): the researcher explains the adaptation
  system (for Experimental participants) or the static system (for
  Control participants) and answers any questions.
\end{enumerate}

Total session duration: approximately 65 minutes per participant.

\paragraph{3.8.5 Analysis Plan}\label{analysis-plan}

\textbf{Quantitative analysis} is implemented in
\texttt{E:\textbackslash{}Projects\textbackslash{}fyp\textbackslash{}v2\textbackslash{}evaluation\textbackslash{}analysis\textbackslash{}stats\_analysis.py}
using the following steps:

\begin{enumerate}
\def\labelenumi{\arabic{enumi}.}
\item
  \textbf{Data preparation:} GEQ raw item scores are imported and
  subscales computed per the scoring key (reverse-scoring items 5, 8, 13
  as required by subscale). NPS-3 and miniPXI scores are computed as
  item means. Behavioral metrics are extracted from the SQLite telemetry
  database.
\item
  \textbf{Normality assessment:} Shapiro-Wilk tests
  (scipy.stats.shapiro) are conducted per DV per condition. With n = 25
  per group, Shapiro-Wilk is the appropriate test (valid for n
  \textless{} 50). If normality is violated (p \textless{} 0.05), the
  Mann-Whitney U test is used instead of Welch's t-test for that DV and
  this deviation is reported.
\item
  \textbf{Primary hypothesis test (H1):} Welch's independent-samples
  t-test (unequal variances, two-tailed) comparing GEQ Flow subscale
  scores between Experimental and Control conditions. Effect size is
  reported as Cohen's d computed via the pingouin library.
\item
  \textbf{Secondary hypothesis tests (H2, H3):} Identical procedure
  applied to GEQ Immersion and NPS-3. Multiple comparison correction
  (Bonferroni) is applied across the three primary hypotheses, yielding
  an adjusted alpha of α = 0.0167 per test.
\item
  \textbf{Reliability:} Cronbach's alpha is computed for each subscale
  using the formula implemented in stats\_analysis.py.
\item
  \textbf{Behavioural metrics:} Welch's t-tests are conducted for
  voluntary\_npc, dialogue\_read\_ratio, and exploration\_coverage as
  exploratory analyses without multiple comparison correction
  (designated as such in the reporting).
\end{enumerate}

\textbf{Qualitative analysis} follows the six-phase Reflexive Thematic
Analysis procedure of Braun and Clarke (2022): (1) familiarization with
data through reading transcripts; (2) generating initial codes across
the full dataset; (3) constructing candidate themes by clustering codes;
(4) reviewing themes against the coded data and the full corpus; (5)
defining and naming themes with clear boundaries; (6) writing up with
representative quotes. The researcher keeps an analytic memo after each
interview to support reflexivity.

\begin{center}\rule{0.5\linewidth}{0.5pt}\end{center}

\subsection{Chapter 4: System
Implementation}\label{chapter-4-system-implementation}

\subsubsection{4.1 Technology Stack}\label{technology-stack}

The system is implemented across two platforms: a Python backend service
and a Unity C\# plugin. Table {[}see below{]} summarizes the key
dependencies.

\textbf{Backend dependencies:}

\begin{longtable}[]{@{}lll@{}}
\toprule\noalign{}
Component & Technology & Version \\
\midrule\noalign{}
\endhead
\bottomrule\noalign{}
\endlastfoot
Runtime & Python & 3.14 \\
API framework & FastAPI & 0.115+ \\
Data validation & Pydantic & v2 \\
Database & SQLite & (via Python stdlib) \\
RL framework & Stable-Baselines3 & 2.x \\
RL environment & Gymnasium & 0.29+ \\
Embeddings & sentence-transformers & latest \\
LLM client & httpx (Ollama API) & 0.27+ \\
Numerical & NumPy & 2.x \\
Statistics & SciPy, pingouin & latest \\
Data analysis & pandas & 2.x \\
\end{longtable}

\textbf{Unity plugin dependencies:}

\begin{longtable}[]{@{}lll@{}}
\toprule\noalign{}
Component & Technology & Version \\
\midrule\noalign{}
\endhead
\bottomrule\noalign{}
\endlastfoot
Engine & Unity & 6 LTS \\
Language & C\# & .NET 9 \\
HTTP client & UnityWebRequest & built-in \\
JSON & Newtonsoft.Json & 13.x \\
UI & TextMeshPro & built-in \\
\end{longtable}

\textbf{Inference:}

\begin{longtable}[]{@{}ll@{}}
\toprule\noalign{}
Component & Technology \\
\midrule\noalign{}
\endhead
\bottomrule\noalign{}
\endlastfoot
LLM runtime & Ollama (local) \\
Primary model & Phi-3.5 Mini (3.8B) \\
Fallback model & Llama 3.2 3B \\
\end{longtable}

\subsubsection{4.2 Backend Service}\label{backend-service}

The FastAPI backend (\texttt{backend/main.py}) exposes four primary
endpoints:

\textbf{Appendix E} documents the complete API specification. The key
endpoints are:

\begin{itemize}
\tightlist
\item
  \texttt{POST\ /telemetry/batch} --- accepts a TelemetryBatch payload,
  persists it to SQLite, triggers player modeling, and initiates the
  adaptation cycle.
\item
  \texttt{GET\ /narrative/generate} --- returns the latest
  NarrativeContent for a given session\_id, blocking until generation is
  complete.
\item
  \texttt{GET\ /player/model/\{player\_id\}} --- returns the current
  PlayerModel for a given player, including Hexad profile, Flow state,
  and session embedding.
\item
  \texttt{WebSocket\ /ws/telemetry/\{session\_id\}} --- low-latency
  telemetry streaming channel for real-time applications.
\end{itemize}

The server is configured with CORS middleware to accept connections from
the Unity client's localhost origin. All data is persisted to a SQLite
database, with separate tables for telemetry batches, player models, and
narrative events. Pydantic v2 models handle all input validation and
output serialization.

\subsubsection{4.3 Database Layer}\label{database-layer}

The SQLite schema uses three primary tables:

\begin{enumerate}
\def\labelenumi{\arabic{enumi}.}
\tightlist
\item
  \textbf{telemetry\_batches} --- stores all 23 TelemetryBatch fields
  plus a server-side received\_at timestamp. Indexed on (player\_id,
  session\_id, window\_start).
\item
  \textbf{player\_models} --- stores PlayerModel snapshots indexed on
  (player\_id, session\_id, timestamp). Each telemetry POST triggers an
  upsert of the computed model.
\item
  \textbf{narrative\_events} --- stores NarrativeContent records
  (action, content\_type, content, speaker, emotional\_tone, fallback
  flag) with foreign keys to session\_id and a generation\_timestamp.
  This log enables post-hoc analysis of which narrative actions were
  fired and at what points in each session.
\end{enumerate}

SQLite was selected over PostgreSQL for simplicity of deployment (no
external service required), study scale (N = 50 generates modest data
volumes), and the single-writer access pattern of the study setup.

\subsubsection{4.4 Player Modeling
Implementation}\label{player-modeling-implementation}

The player modeling pipeline is implemented as three sequential classes
called within the same request handler:

\begin{verbatim}
TelemetryBatch[] (from SQLite)
    → FeatureExtractor.extract()  → dict[str, float]  (11 features)
    → HexadProfiler.compute()     → HexadProfile       (6 dims)
    → FlowClassifier.classify()   → (FlowState, float) (state + confidence)
    → PlayerModel (assembled)
\end{verbatim}

The FeatureExtractor aggregates all TelemetryBatch records for the
current session from SQLite before computing features, ensuring that
features reflect the full session history rather than only the most
recent window. This cumulative computation is appropriate for Hexad
profiling (which requires sufficient behavioral evidence) while the
FlowClassifier uses only the most recent window's features to classify
the instantaneous state.

A notable implementation detail is the denominator protection pattern
applied throughout: \texttt{max(deaths,\ 1)},
\texttt{max(obj\_attempted,\ 1)}, etc. These guards prevent division by
zero in the opening minutes of a session when many counters are still
zero, and are implemented uniformly across both FeatureExtractor and
HexadProfiler.

\subsubsection{4.5 Narrative Pipeline
Implementation}\label{narrative-pipeline-implementation}

The narrative pipeline is invoked asynchronously following each player
model update:

\begin{Shaded}
\begin{Highlighting}[]
\ControlFlowTok{async} \KeywordTok{def}\NormalTok{ generate\_narrative(session\_id: }\BuiltInTok{str}\NormalTok{, player\_model: PlayerModel) }\OperatorTok{{-}\textgreater{}}\NormalTok{ NarrativeContent:}
\NormalTok{    action }\OperatorTok{=}\NormalTok{ agent.select\_action(player\_model)}
\NormalTok{    query }\OperatorTok{=} \SpecialStringTok{f"}\SpecialCharTok{\{}\NormalTok{player\_model}\SpecialCharTok{.}\NormalTok{flow\_state}\SpecialCharTok{.}\NormalTok{value}\SpecialCharTok{\}}\SpecialStringTok{ }\SpecialCharTok{\{}\NormalTok{player\_model}\SpecialCharTok{.}\NormalTok{current\_state}\SpecialCharTok{.}\NormalTok{value}\SpecialCharTok{\}}\SpecialStringTok{ }\SpecialCharTok{\{}\NormalTok{player\_model}\SpecialCharTok{.}\NormalTok{current\_location}\SpecialCharTok{\}}\SpecialStringTok{"}
\NormalTok{    lore\_chunks }\OperatorTok{=}\NormalTok{ retriever.retrieve(query, top\_k}\OperatorTok{=}\DecValTok{3}\NormalTok{)}
\NormalTok{    lore\_texts }\OperatorTok{=}\NormalTok{ [c[}\StringTok{"text"}\NormalTok{] }\ControlFlowTok{for}\NormalTok{ c }\KeywordTok{in}\NormalTok{ lore\_chunks]}
\NormalTok{    memory\_ctx }\OperatorTok{=}\NormalTok{ memory.get\_context\_summary()}
\NormalTok{    system\_p, user\_p }\OperatorTok{=}\NormalTok{ prompt\_builder.build(}
\NormalTok{        player\_model, action,}
\NormalTok{        location}\OperatorTok{=}\NormalTok{player\_model.current\_location,}
\NormalTok{        quest\_stage}\OperatorTok{=}\NormalTok{player\_model.current\_objective,}
\NormalTok{        lore\_context}\OperatorTok{=}\NormalTok{lore\_texts,}
\NormalTok{        memory\_context}\OperatorTok{=}\NormalTok{memory\_ctx,}
\NormalTok{    )}
\NormalTok{    content }\OperatorTok{=} \ControlFlowTok{await}\NormalTok{ ollama.generate(system\_p, user\_p, action)}
    \ControlFlowTok{if} \KeywordTok{not}\NormalTok{ content.fallback:}
\NormalTok{        memory.add\_event(content.content, importance}\OperatorTok{=}\FloatTok{0.7}\NormalTok{)}
    \ControlFlowTok{return}\NormalTok{ content}
\end{Highlighting}
\end{Shaded}

The RAG query string is constructed from the Flow state, current
activity, and location to maximize retrieval relevance: a player in
ANXIETY during IN\_COMBAT in the Dungeon will retrieve lore chunks about
the dungeon's dangers, providing contextually grounded guidance. The
importance score of 0.7 assigned to new memory events reflects a
baseline estimate; a more sophisticated implementation could derive
importance from the narrative action type (e.g., LORE\_REWARD events
receive higher importance than NO\_CHANGE events).

\subsubsection{4.6 RL Adaptation
Implementation}\label{rl-adaptation-implementation}

The NarrativeAdaptationEnv is registered as a standard Gymnasium
environment and trained using Stable-Baselines3 PPO with the
hyperparameters in Table 9.

\textbf{Table 9 --- PPO hyperparameters}

\begin{longtable}[]{@{}ll@{}}
\toprule\noalign{}
Hyperparameter & Value \\
\midrule\noalign{}
\endhead
\bottomrule\noalign{}
\endlastfoot
Policy & MlpPolicy (2 x 64 hidden layers) \\
Learning rate & 3 × 10⁻⁴ \\
n\_steps & 512 \\
batch\_size & 64 \\
n\_epochs & 10 \\
gamma (γ) & 0.95 \\
gae\_lambda & 0.95 \\
clip\_range & 0.2 \\
n\_envs (vectorized) & 4 \\
total\_timesteps & 200,000 \\
\end{longtable}

Training uses four parallel vectorized environments
(\texttt{make\_vec\_env}) to improve sample efficiency. The trained
model is saved to \texttt{./data/ppo\_narrative\_agent.zip} and loaded
on subsequent server starts, avoiding the 200,000-step training overhead
at inference time.

The NarrativeAgent.select\_action() method constructs the 7-dimensional
observation vector from the PlayerModel (flow\_score, normalized
challenge\_skill\_ratio, explorer, socializer, achiever, disruptor,
session\_progress) and calls model.predict(obs, deterministic=False).
Stochastic prediction is used at inference time to maintain behavioral
diversity; deterministic=True would always select the
maximum-probability action, potentially producing repetitive narrative
patterns.

\subsubsection{4.7 Unity Plugin
Implementation}\label{unity-plugin-implementation}

The Unity plugin is structured as a package-compatible set of C\#
scripts under \texttt{Assets/NarrativePersonalization/}. The
PlayerDataLogger accumulates telemetry via direct field increments in
Update() and fires the HTTP POST at fixed 5-second intervals using a
Coroutine.

The NarrativeManager uses a polling architecture: immediately following
each telemetry POST's response, it fires a GET request to
/narrative/generate. The 8-second Ollama generation timeout is handled
by displaying no new content rather than blocking the game loop. The
ContentInjector implements a state machine with three states: IDLE,
DISPLAYING, and COOLING\_DOWN. Content queued during COOLING\_DOWN is
held and injected after the 3-second minimum spacing, ensuring narrative
delivery does not overwhelm the player during intense combat sequences.

\subsubsection{4.8 Testing and Validation}\label{testing-and-validation}

\textbf{Unit tests} were written for the FeatureExtractor (edge cases:
empty batch list, zero-denominator fields, maximum value clamping),
HexadProfiler (initialization with empty batches returns 0.5 defaults),
FlowClassifier (boundary conditions at ratio 0.75 and 1.30, APATHY
compound condition), and compute\_reward (all state transitions, streak
bonus triggering, repetition penalty).

\textbf{Integration tests} validated the end-to-end pipeline from a mock
TelemetryBatch POST to a returned NarrativeContent JSON, confirming that
Flow state correctly influences the generated tone and that the fallback
mechanism activates when Ollama is unavailable.

\textbf{Validation of the simulated MDP} was conducted by verifying that
the trained PPO agent converged to a mean episodic reward substantially
above zero (indicating that the agent learned to prefer helpful actions)
and that the agent's action distribution in ANXIETY states was dominated
by PROVIDE\_GUIDANCE and LOWER\_STAKES --- consistent with the heuristic
--- rather than random.

\textbf{Lore retrieval validation} was conducted by manually querying
the retriever with ten representative in-session queries and verifying
that retrieved chunks were semantically relevant. All ten manual queries
returned at least one clearly relevant chunk in the top-3 results.

\begin{center}\rule{0.5\linewidth}{0.5pt}\end{center}

\subsection{Chapter 5: Results}\label{chapter-5-results}

\begin{quote}
\textbf{Note:} This chapter presents structured result templates. All
{[}INSERT: \ldots{]} placeholders will be replaced with data following
completion of the user study. All table cells, statistical values, and
quotations marked as {[}INSERT{]} are placeholders and do not represent
actual findings.
\end{quote}

\subsubsection{5.1 Participant
Demographics}\label{participant-demographics}

A total of N = 50 participants completed the study (25 Experimental, 25
Control). {[}INSERT: one participant was excluded from analysis due
to\ldots{]} yielding a final sample of {[}INSERT: N{]}. Table 11
summarizes demographics.

\textbf{Table 11 --- Participant demographics {[}PLACEHOLDER{]}}

\begin{longtable}[]{@{}
  >{\raggedright\arraybackslash}p{(\columnwidth - 6\tabcolsep) * \real{0.2500}}
  >{\raggedright\arraybackslash}p{(\columnwidth - 6\tabcolsep) * \real{0.2500}}
  >{\raggedright\arraybackslash}p{(\columnwidth - 6\tabcolsep) * \real{0.2500}}
  >{\raggedright\arraybackslash}p{(\columnwidth - 6\tabcolsep) * \real{0.2500}}@{}}
\toprule\noalign{}
\begin{minipage}[b]{\linewidth}\raggedright
Variable
\end{minipage} & \begin{minipage}[b]{\linewidth}\raggedright
Experimental (n=25)
\end{minipage} & \begin{minipage}[b]{\linewidth}\raggedright
Control (n=25)
\end{minipage} & \begin{minipage}[b]{\linewidth}\raggedright
Total (N=50)
\end{minipage} \\
\midrule\noalign{}
\endhead
\bottomrule\noalign{}
\endlastfoot
Age: Mean (SD) & {[}INSERT{]} & {[}INSERT{]} & {[}INSERT{]} \\
Gender: Female / Male / Non-binary / Prefer not to say & {[}INSERT{]} &
{[}INSERT{]} & {[}INSERT{]} \\
Gaming hours/week: Mean (SD) & {[}INSERT{]} & {[}INSERT{]} &
{[}INSERT{]} \\
RPG experience (1-5): Mean (SD) & {[}INSERT{]} & {[}INSERT{]} &
{[}INSERT{]} \\
\end{longtable}

No significant baseline differences between conditions were observed on
age (t({[}INSERT{]}) = {[}INSERT{]}, p = {[}INSERT{]}) or gaming hours
(t({[}INSERT{]}) = {[}INSERT{]}, p = {[}INSERT{]}), supporting the
assumption of pre-intervention equivalence.

\subsubsection{5.2 Quantitative Results}\label{quantitative-results}

\paragraph{5.2.1 GEQ Flow Subscale (H1)}\label{geq-flow-subscale-h1}

\textbf{Table 12a --- GEQ Flow subscale {[}PLACEHOLDER{]}}

\begin{longtable}[]{@{}lllll@{}}
\toprule\noalign{}
Condition & n & M & SD & 95\% CI \\
\midrule\noalign{}
\endhead
\bottomrule\noalign{}
\endlastfoot
Experimental & 25 & {[}INSERT{]} & {[}INSERT{]} & {[}INSERT,
INSERT{]} \\
Control & 25 & {[}INSERT{]} & {[}INSERT{]} & {[}INSERT, INSERT{]} \\
\end{longtable}

Welch's independent-samples t-test: t({[}INSERT{]}) = {[}INSERT{]}, p =
{[}INSERT{]}, Cohen's d = {[}INSERT{]}.

{[}INSERT: H1 was / was not supported. The Experimental condition scored
{[}direction{]} higher on the GEQ Flow subscale (M = {[}INSERT{]}, SD =
{[}INSERT{]}) compared to the Control condition (M = {[}INSERT{]}, SD =
{[}INSERT{]}). The effect size was {[}magnitude descriptor{]} (d =
{[}INSERT{]}).{]}

Shapiro-Wilk normality test: Experimental W = {[}INSERT{]}, p =
{[}INSERT{]}; Control W = {[}INSERT{]}, p = {[}INSERT{]}. {[}INSERT:
normality was / was not violated; if violated, Mann-Whitney U was
substituted.{]}

Cronbach's alpha for the GEQ Flow subscale: α = {[}INSERT{]}.

\paragraph{5.2.2 GEQ Immersion Subscale
(H2)}\label{geq-immersion-subscale-h2}

\textbf{Table 12b --- GEQ Immersion subscale {[}PLACEHOLDER{]}}

\begin{longtable}[]{@{}lllll@{}}
\toprule\noalign{}
Condition & n & M & SD & 95\% CI \\
\midrule\noalign{}
\endhead
\bottomrule\noalign{}
\endlastfoot
Experimental & 25 & {[}INSERT{]} & {[}INSERT{]} & {[}INSERT,
INSERT{]} \\
Control & 25 & {[}INSERT{]} & {[}INSERT{]} & {[}INSERT, INSERT{]} \\
\end{longtable}

Welch's t-test: t({[}INSERT{]}) = {[}INSERT{]}, p = {[}INSERT{]},
Cohen's d = {[}INSERT{]}.

{[}INSERT: H2 was / was not supported. Cronbach's alpha for GEQ
Immersion: α = {[}INSERT{]}.{]}

\paragraph{5.2.3 Personalization Perception
(H3)}\label{personalization-perception-h3}

\textbf{Table 12c --- NPS-3 Personalization Scale {[}PLACEHOLDER{]}}

\begin{longtable}[]{@{}lllll@{}}
\toprule\noalign{}
Condition & n & M & SD & 95\% CI \\
\midrule\noalign{}
\endhead
\bottomrule\noalign{}
\endlastfoot
Experimental & 25 & {[}INSERT{]} & {[}INSERT{]} & {[}INSERT,
INSERT{]} \\
Control & 25 & {[}INSERT{]} & {[}INSERT{]} & {[}INSERT, INSERT{]} \\
\end{longtable}

Welch's t-test: t({[}INSERT{]}) = {[}INSERT{]}, p = {[}INSERT{]},
Cohen's d = {[}INSERT{]}.

{[}INSERT: H3 was / was not supported. NPS-3 Cronbach's alpha: α =
{[}INSERT{]}. Item-level means: Item 1 = {[}INSERT{]}, Item 2 =
{[}INSERT{]}, Item 3 = {[}INSERT{]}.{]}

\paragraph{5.2.4 Behavioural Metrics}\label{behavioural-metrics}

\textbf{Table 12d --- Behavioural metrics {[}PLACEHOLDER{]}}

\begin{longtable}[]{@{}
  >{\raggedright\arraybackslash}p{(\columnwidth - 10\tabcolsep) * \real{0.1667}}
  >{\raggedright\arraybackslash}p{(\columnwidth - 10\tabcolsep) * \real{0.1667}}
  >{\raggedright\arraybackslash}p{(\columnwidth - 10\tabcolsep) * \real{0.1667}}
  >{\raggedright\arraybackslash}p{(\columnwidth - 10\tabcolsep) * \real{0.1667}}
  >{\raggedright\arraybackslash}p{(\columnwidth - 10\tabcolsep) * \real{0.1667}}
  >{\raggedright\arraybackslash}p{(\columnwidth - 10\tabcolsep) * \real{0.1667}}@{}}
\toprule\noalign{}
\begin{minipage}[b]{\linewidth}\raggedright
Metric
\end{minipage} & \begin{minipage}[b]{\linewidth}\raggedright
Experimental M (SD)
\end{minipage} & \begin{minipage}[b]{\linewidth}\raggedright
Control M (SD)
\end{minipage} & \begin{minipage}[b]{\linewidth}\raggedright
t
\end{minipage} & \begin{minipage}[b]{\linewidth}\raggedright
p
\end{minipage} & \begin{minipage}[b]{\linewidth}\raggedright
d
\end{minipage} \\
\midrule\noalign{}
\endhead
\bottomrule\noalign{}
\endlastfoot
Voluntary NPC interactions & {[}INSERT{]} & {[}INSERT{]} & {[}INSERT{]}
& {[}INSERT{]} & {[}INSERT{]} \\
Dialogue read ratio & {[}INSERT{]} & {[}INSERT{]} & {[}INSERT{]} &
{[}INSERT{]} & {[}INSERT{]} \\
Exploration coverage & {[}INSERT{]} & {[}INSERT{]} & {[}INSERT{]} &
{[}INSERT{]} & {[}INSERT{]} \\
miniPXI total score & {[}INSERT{]} & {[}INSERT{]} & {[}INSERT{]} &
{[}INSERT{]} & {[}INSERT{]} \\
\end{longtable}

{[}INSERT: Exploratory analysis indicated that participants in the
Experimental condition {[}direction{]} more/fewer voluntary NPC
interactions (M = {[}INSERT{]}) compared to Control (M = {[}INSERT{]}),
suggesting {[}interpretation{]}.{]}

\subsubsection{5.3 Qualitative Results (Thematic
Analysis)}\label{qualitative-results-thematic-analysis}

Semi-structured interviews were conducted with a subset of {[}INSERT:
n{]} participants ({[}INSERT: n{]} Experimental, {[}INSERT: n{]}
Control) across sessions lasting {[}INSERT: M{]} minutes on average.
Following the six-phase Reflexive Thematic Analysis procedure (Braun \&
Clarke, 2022), {[}INSERT: n{]} initial codes were generated,
subsequently organized into {[}INSERT: n{]} candidate themes, and
refined into {[}INSERT: n{]} final themes. Table 13 presents the themes
with representative participant quotations.

\textbf{Table 13 --- Qualitative themes and representative quotations
{[}PLACEHOLDER{]}}

\begin{longtable}[]{@{}
  >{\raggedright\arraybackslash}p{(\columnwidth - 4\tabcolsep) * \real{0.3333}}
  >{\raggedright\arraybackslash}p{(\columnwidth - 4\tabcolsep) * \real{0.3333}}
  >{\raggedright\arraybackslash}p{(\columnwidth - 4\tabcolsep) * \real{0.3333}}@{}}
\toprule\noalign{}
\begin{minipage}[b]{\linewidth}\raggedright
Theme
\end{minipage} & \begin{minipage}[b]{\linewidth}\raggedright
Description
\end{minipage} & \begin{minipage}[b]{\linewidth}\raggedright
Representative Quote
\end{minipage} \\
\midrule\noalign{}
\endhead
\bottomrule\noalign{}
\endlastfoot
{[}INSERT: Theme 1 name{]} & {[}INSERT: description{]} & {[}INSERT:
``participant quote'' (Exp/Ctrl, P{[}n{]}){]} \\
{[}INSERT: Theme 2 name{]} & {[}INSERT: description{]} & {[}INSERT:
``participant quote'' (Exp/Ctrl, P{[}n{]}){]} \\
{[}INSERT: Theme 3 name{]} & {[}INSERT: description{]} & {[}INSERT:
``participant quote'' (Exp/Ctrl, P{[}n{]}){]} \\
{[}INSERT: Theme 4 name{]} & {[}INSERT: description{]} & {[}INSERT:
``participant quote'' (Exp/Ctrl, P{[}n{]}){]} \\
\end{longtable}

{[}INSERT: Anticipated themes based on the literature and system design
include: (1) narrative responsiveness --- participants in the
Experimental condition describing the story as ``reactive'' or ``aware''
of their play style; (2) immersive engagement --- descriptions of being
``pulled into'' the world at moments of narrative injection; (3)
artificiality awareness --- participants noticing inconsistencies or
repetitive patterns in AI-generated content; (4) emotional resonance ---
the emotional tone of narrative injections matching or mismatching the
player's affective state.{]}

\subsubsection{5.4 Summary of Findings}\label{summary-of-findings}

{[}INSERT: A narrative summary of the overall pattern of results across
quantitative and qualitative strands will be provided here following
data collection. The summary will address: whether the primary
hypothesis (H1, GEQ Flow) was supported; whether the secondary
hypotheses (H2, H3) were supported; what the behavioural metrics
indicate about engagement differences; and the dominant qualitative
themes emerging from the interview data.{]}

\begin{center}\rule{0.5\linewidth}{0.5pt}\end{center}

\subsection{Chapter 6: Discussion}\label{chapter-6-discussion}

\subsubsection{6.1 Interpretation of Quantitative
Findings}\label{interpretation-of-quantitative-findings}

The central prediction of the research --- that PPO-driven narrative
adaptation would produce measurably higher Flow experience --- rests on
the theoretical chain from the adaptation engine through the player
modeling module to the GEQ instrument. If the PPO agent successfully
learns to apply narratively appropriate actions in response to Flow
state signals, and if those narrative interventions succeed in shifting
the player's psychological state toward FLOW, then the accumulated Flow
experience across the session should register as a higher GEQ Flow
subscale score.

{[}INSERT: Interpret the actual results when available. Example
structures follow:{]}

\textbf{If H1 is supported:} The significant difference in GEQ Flow
scores (d = {[}INSERT{]}) supports the core theoretical claim that
real-time narrative personalization can improve Flow experience. The
effect size is {[}small/medium/large{]} by Cohen's (1988) convention,
suggesting that narrative adaptation is a meaningful lever for
engagement even within a relatively short (30-minute) session. The
convergent evidence from the behavioural metrics --- specifically the
{[}higher/lower{]} dialogue read ratio in the Experimental condition ---
suggests that players were not merely reporting higher Flow
retrospectively but were also demonstrably more engaged with the
narrative content during play.

\textbf{If H1 is not supported:} The absence of a significant difference
in GEQ Flow scores challenges the assumption that within-session
narrative adaptation produces observable Flow differences within a
30-minute session. Several explanations merit consideration. First, the
PPO agent was trained on a simulated environment whose transition
dynamics only approximately match real player behavior; policy transfer
to the real game environment may be limited. Second, the 30-minute
session may be insufficient for the player modeling module to accumulate
enough telemetry to produce reliable Hexad profiles and accurate Flow
state classifications; narrative adaptation in the final third of a
session could not compensate for the earlier period of profile-building.
Third, the effect of narrative tone on Flow may be smaller than
anticipated relative to the contribution of mechanical game design
factors (challenge level, control responsiveness).

{[}INSERT: Discuss H2 (Immersion) and H3 (Personalization) results in
relation to each other and to H1. Note any divergence between hypothesis
outcomes --- e.g., if H3 is supported but H1 is not, this suggests that
players perceived the narrative as personalized without this perception
translating to measurable Flow improvement, which would be theoretically
interesting.{]}

\subsubsection{6.2 Themes from Interviews}\label{themes-from-interviews}

{[}INSERT: Discuss the qualitative themes in relation to the
quantitative results and the existing literature. Address convergence
and divergence between the two strands. Example discussion points:{]}

The interview data are expected to illuminate the mechanisms by which
narrative adaptation influenced (or failed to influence) experience. A
theme of \textbf{narrative responsiveness} --- players describing the
story as attuned to their behavior --- would provide direct qualitative
support for the personalization hypothesis (H3) and would help explain
any quantitative differences in NPS-3 scores. Conversely, a theme of
\textbf{artificiality awareness} --- players noticing repetitive
patterns, inconsistent tone, or implausible content --- would identify a
key failure mode of the LLM generation pipeline and suggest directions
for improvement (e.g., larger context windows, more diverse fallback
content, human-in-the-loop content review).

The Flow-related interview questions (Section 3 of the interview guide)
are expected to reveal whether narrative injections were perceived as
disruptive or supportive of engagement. A recurring finding in studies
of game storytelling is that narrative interruptions during
high-engagement moments (e.g., entering a combat encounter) are
perceived negatively even when the content is high quality. The
ContentInjector's 3-second spacing mechanism was designed to mitigate
this, but interview data will reveal whether it was sufficient.

{[}INSERT: Link themes to the specific interview guide questions
(Appendix D) and to specific moments in the game sessions as identified
through triangulation with the telemetry data.{]}

\subsubsection{6.3 Novelty and
Contributions}\label{novelty-and-contributions}

This work makes the following original contributions to the fields of
game narrative, player modeling, and applied reinforcement learning:

\textbf{Contribution 1: Telemetry-derived Hexad profiling.} Prior work
on Hexad-based game adaptation has relied exclusively on pre-play
questionnaire scores (Tondello et al., 2016). The present system
computes six-dimensional continuous Hexad profiles from behavioral
telemetry, enabling within-session updating without any player
self-report. The mapping from telemetry signals to Hexad dimensions
(Table 4) provides a reusable operational template for game developers
seeking to implement Hexad-adaptive systems.

\textbf{Contribution 2: RAG-grounded dynamic narrative.} While PANGeA
(Todd et al., 2024) demonstrated LLM-based procedural narrative grounded
in prompt-delivered context, the present work extends this to a dynamic
retrieval paradigm in which the specific lore passages delivered to the
LLM are selected at runtime based on the player's current context. This
ensures that generated narrative is semantically relevant to the
player's location and activity, not merely to the game world in general.

\textbf{Contribution 3: Episodic session memory in the LLM prompt.} The
implementation of importance-weighted episodic memory inspired by EM-LLM
(Fountas et al., 2024) enables narrative continuity across a play
session. To the author's knowledge, this is the first application of
episodic memory compression to real-time game narrative generation. The
mechanism ensures that generated content can reference prior events
without exceeding the LLM's context window.

\textbf{Contribution 4: PPO for narrative action selection.}
Reinforcement learning has been applied extensively to mechanical game
adaptation (DDA) but not, to the author's knowledge, to narrative action
selection. Framing narrative tone selection as an RL problem --- with a
learned policy conditioned on player psychological state --- represents
a meaningful conceptual advance over rule-based narrative adaptation
systems.

\textbf{Contribution 5: Flow-based MDP reward.} The reward function
directly operationalizes Csikszentmihalyi's (1990) challenge-skill
balance theory, providing a principled theoretical grounding for the
adaptation objective. The shaped reward (streak bonus, transition bonus,
repetition penalty) encourages the learned policy to actively maintain
Flow rather than merely avoiding non-Flow states.

\textbf{Contribution 6: Complete closed-loop integration.} The
five-module architecture is novel in its end-to-end integration: from
raw telemetry capture in Unity, through behavioral feature extraction
and player typing, through RL action selection, through RAG-grounded LLM
generation, to in-game content injection --- all within a single
deployable plugin. No prior system in the literature integrates all of
these components.

\subsubsection{6.4 Limitations}\label{limitations}

Several limitations of the present work merit acknowledgment.

\textbf{Simulation fidelity.} The PPO agent is trained on a simulated
environment with hand-specified transition probabilities (p = 0.75 for
successful helpful actions). Real player behavior is substantially more
complex and variable; the simulation's Markovian assumption (that
transitions depend only on current state and action) is almost certainly
violated in practice, as player engagement trajectories have temporal
dependencies extending well beyond the current state. Policy transfer
from simulation to real gameplay may therefore be imperfect.

\textbf{Rule-based Hexad profiling.} The mapping from telemetry signals
to Hexad dimensions (Table 4) is based on theoretical alignment rather
than empirical validation against questionnaire-derived ground truth.
Players may exhibit behavioral patterns that are systematically
misclassified --- for instance, a player who moves slowly due to
hardware input lag rather than deliberate exploration will register high
explore\_rate and be profiled as an Explorer. Future work should
validate the telemetry-derived scores against matched questionnaire
data.

\textbf{LLM generation quality.} Phi-3.5 Mini (3.8B parameters) is a
small language model constrained by its capacity relative to larger
frontier models. While JSON-mode generation and the 8-second timeout
mitigate failures, the quality of generated narrative --- its fluency,
creativity, and character voice consistency --- is limited by the model
size. Some generated outputs may be repetitive, tonally inconsistent, or
fail to maintain character voice despite the explicit voice descriptions
in the prompt.

\textbf{Single-session study.} The study captures a single 30-minute
session per participant. The long-term effects of narrative
personalization --- whether players develop a richer relationship with
the game world over multiple sessions, whether the Hexad profiling
becomes more accurate with more data --- cannot be assessed within this
design.

\textbf{Testbed ecological validity.} The testbed game, while designed
to be engaging, is a simplified prototype with three scenes and three
NPCs. Findings may not generalize to commercially developed games with
more complex narrative architectures, greater visual fidelity, or more
sophisticated combat systems.

\textbf{Sample size.} With n = 25 per condition, the study is powered to
detect medium effects (d ≈ 0.5) but not small effects. If the true
effect of narrative personalization is small, the study will be
underpowered to detect it. The convenience sample from a university
population also limits demographic diversity.

\subsubsection{6.5 Threats to Validity}\label{threats-to-validity}

\textbf{Internal validity.} The primary threat is the potential for
demand characteristics: participants who infer (correctly or
incorrectly) that they are in the experimental condition may respond
more favorably to questionnaire items about personalization and flow.
The blind study design (participants are not told their condition during
the session) mitigates but does not eliminate this threat, as observant
participants may notice that the narrative seems responsive to their
behavior. The debrief is conducted after questionnaire completion to
prevent retrospective contamination.

\textbf{Construct validity.} The GEQ Flow subscale (items 5, 13, 25, 28,
31) measures the phenomenological experience of flow as reported
immediately post-session, not flow during gameplay. Retrospective recall
of engagement states is subject to recency bias and mood-congruent
recall effects. The absence of a physiological measure (e.g., galvanic
skin response, heart rate variability) limits confidence that the GEQ
score reflects the within-session psychological state targeted by the
adaptation engine.

\textbf{External validity.} As noted in the limitations section, the
testbed game and the student convenience sample limit generalizability.
The findings should be treated as applicable to RPG-adjacent games with
narrative content and to a young, gaming-experienced population until
replicated in broader contexts.

\textbf{Statistical conclusion validity.} With N = 50, the study is not
robustly powered for small effect sizes, and the Bonferroni correction
for multiple comparisons (three hypotheses) reduces power further. The
choice of Welch's t-test (rather than the standard Student's t-test)
provides robustness against unequal variances but assumes approximate
normality. If normality is substantially violated, the Mann-Whitney U
substitution is less powerful and interpretable as a test of
distributional differences rather than means.

\begin{center}\rule{0.5\linewidth}{0.5pt}\end{center}

\subsection{Chapter 7: Conclusion}\label{chapter-7-conclusion}

\subsubsection{7.1 Summary of
Contributions}\label{summary-of-contributions}

This dissertation has presented a complete design, implementation, and
evaluation methodology for an AI-driven closed-loop narrative
personalization plugin for Unity games. The system integrates five
modules --- telemetry collection, player modeling, narrative generation,
RL adaptation, and content injection --- into a coherent pipeline that
operates within 5-second adaptation cycles.

The six principal contributions are: (1) a method for deriving
continuous Hexad player type profiles from behavioral telemetry without
requiring a player questionnaire; (2) a RAG-grounded narrative
generation pipeline using sentence-transformers cosine retrieval over a
structured lore corpus; (3) an importance-weighted episodic session
memory that provides LLM narrative continuity across a play session; (4)
a PPO reinforcement learning agent trained to select among eight
narrative tone actions based on player psychological state; (5) a
Flow-based MDP reward function operationalizing Csikszentmihalyi's
challenge-skill balance theory; and (6) a complete closed-loop
integration of all five components within a deployable Unity plugin.

The evaluation design --- a between-subjects study (N = 50) using the
GEQ, miniPXI, and custom NPS-3 scale, supplemented by reflexive thematic
analysis of post-play interviews --- provides a rigorous mixed-methods
framework for assessing the system's impact on Flow experience,
immersion, and perceived personalization. The statistical analysis
pipeline is implemented as a reusable Python script
(\texttt{evaluation/analysis/stats\_analysis.py}) that will produce
APA-formatted results tables directly from the collected data.

\subsubsection{7.2 Future Work}\label{future-work}

Several directions for future work present themselves.

\textbf{Empirical validation of telemetry-derived Hexad profiling.} The
most immediate extension is a validation study comparing
telemetry-derived Hexad scores against matched questionnaire scores (the
Gamification User Types Hexad Scale; Tondello et al., 2016) from the
same participants. This would quantify the construct validity of the
behavioral proxy and identify which dimensions are most accurately
captured by the current feature-to-dimension mappings.

\textbf{Online RL adaptation.} The current PPO agent is trained offline
on a simulated environment. An online or hybrid approach --- in which
the agent continues to update its policy based on real player
interaction data collected during study --- would produce a policy
increasingly tailored to the real game environment. Careful management
of exploration-exploitation tradeoffs and safety constraints (preventing
the policy from selecting narratively incoherent actions during
exploration) would be necessary.

\textbf{Larger LLMs and fine-tuning.} Replacing Phi-3.5 Mini with a
larger model (e.g., Llama 3.1 8B or a fine-tuned narrative generation
model) could substantially improve generation quality, character voice
consistency, and tonal nuance. Fine-tuning on a corpus of high-quality
dark fantasy dialogue --- potentially generated by a larger model with
human curation --- would align the model's prior distribution more
closely with the game world's stylistic requirements.

\textbf{Multi-session study.} A longitudinal study spanning three to
five sessions per participant would reveal whether the personalization
system produces cumulative benefits: whether the Hexad profile becomes
more accurate and the adaptation policy more appropriate as more session
data accumulates, and whether players develop a richer relationship with
the narrative over repeated play.

\textbf{Evaluation in a commercial game context.} Collaborating with an
independent game developer to deploy the plugin in an in-development
commercial title would provide a more ecologically valid test and access
to a more diverse player population than a convenience sample of
students.

\textbf{Voice and audio integration.} The current system delivers
narrative as text. Integrating text-to-speech generation (potentially
with character-specific voice parameters) would create a richer
multimodal experience and align the audio delivery with the generated
emotional tone.

\textbf{Explainability and player agency.} A transparency mode in which
players can see why a particular narrative action was chosen --- ``The
game noticed you seem to be having a difficult time and offered some
guidance'' --- could be studied as a mechanism for building player trust
in the system and supporting meta-cognitive engagement with the
personalization process.

\subsubsection{7.3 Closing Remarks}\label{closing-remarks}

The aspiration behind this research is modest in scope but substantive
in ambition: to demonstrate that the tools now available to AI
researchers --- large language models, reinforcement learning, semantic
retrieval --- can be composed into a practical system that makes games
more attentive to the people who play them. The challenge is not merely
technical but also philosophical: what does it mean for a game's story
to understand a player? What behavioral signals are worth attending to,
and what should be ignored? When does adaptation feel like
personalization and when does it feel like surveillance?

These questions are ultimately design and ethics questions as much as
they are engineering questions, and the present work represents only an
early step toward answering them. The closed-loop system described here
is a proof of concept: a demonstration that the integration is feasible,
that the components can communicate, and that the adaptation cycle can
run within the latency constraints of real-time gameplay. Whether it
produces experiences that players find genuinely enriched by the
adaptation --- whether the personalization is felt rather than merely
inferred --- is an empirical question that the evaluation study will
begin to answer.

The long arc of games as a medium bends toward greater responsiveness to
individual players. Dynamic narrative personalization, grounded in
psychological models and driven by learned adaptive policies, is one
promising path along that arc.

\begin{center}\rule{0.5\linewidth}{0.5pt}\end{center}

\subsection{References}\label{references}

Bartle, R. (1996). Hearts, clubs, diamonds, spades: Players who suit
MUDs. \emph{Journal of MUD Research, 1}(1).
https://mud.co.uk/richard/hcds.htm

Braun, V., \& Clarke, V. (2022). \emph{Thematic analysis: A practical
guide}. SAGE Publications.

Chen, J. (2007). Flow in games (and everything else).
\emph{Communications of the ACM, 50}(4), 31--34.
https://doi.org/10.1145/1232743.1232769

Compton, K., Kybartas, B., \& Mateas, M. (2015). Tracery: An
author-focused generative text tool. In \emph{Interactive Storytelling:
8th International Conference on Interactive Digital Storytelling}
(pp.~154--161). Springer.

Csikszentmihalyi, M. (1990). \emph{Flow: The psychology of optimal
experience}. Harper \& Row.

Faul, F., Erdfelder, E., Lang, A.-G., \& Buchner, A. (2007).
G\emph{Power 3: A flexible statistical power analysis program for the
social, behavioral, and biomedical sciences. }Behavior Research Methods,
39*(2), 175--191. https://doi.org/10.3758/BF03193146

Fountas, Z., Sylaidi, A., Bhatt, U., \& Bhatt, P. (2024).
\emph{Human-like episodic memory for infinite context LLMs}. arXiv
preprint arXiv:2407.09450. {[}Accepted at ICLR 2025{]}

Hunicke, R., \& Chapman, V. (2004). AI for dynamic difficulty adjustment
in games. In \emph{Proceedings of the AAAI Workshop on Challenges in
Game AI} (pp.~91--96). AAAI Press.

IJsselsteijn, W. A., de Kort, Y. A. W., \& Poels, K. (2013). \emph{The
Game Experience Questionnaire}. Eindhoven: Technische Universiteit
Eindhoven.

Kreminski, M., Samuel, B., Melcer, E., \& Wardrip-Fruin, N. (2019).
Evaluating AI-based games through the lens of interactivity. In
\emph{Proceedings of the AAAI Conference on Artificial Intelligence and
Interactive Digital Entertainment (AIIDE)} (pp.~59--65). AAAI Press.

Lewis, P., Perez, E., Piktus, A., Petroni, F., Karpukhin, V., Goyal, N.,
Küttler, H., Lewis, M., Yih, W., Rocktäschel, T., Riedel, S., \& Kiela,
D. (2020). Retrieval-augmented generation for knowledge-intensive NLP
tasks. \emph{Advances in Neural Information Processing Systems
(NeurIPS), 33}, 9459--9474.

Mateas, M., \& Stern, A. (2003). Façade: An experiment in building a
fully-realized interactive drama. In \emph{Proceedings of the Game
Developers Conference (GDC)}. CMP Media.

Microsoft. (2024). \emph{Phi-3.5 Mini Instruct model card}. Microsoft
Research. https://huggingface.co/microsoft/Phi-3.5-mini-instruct

Raffin, A., Hill, A., Gleave, A., Kanervisto, A., Ernestus, M., \&
Dormann, N. (2021). Stable-Baselines3: Reliable reinforcement learning
implementations. \emph{Journal of Machine Learning Research, 22}(268),
1--8. https://jmlr.org/papers/v22/20-1364.html

Ryan, R. M., \& Deci, E. L. (2000). Self-determination theory and the
facilitation of intrinsic motivation, social development, and
well-being. \emph{American Psychologist, 55}(1), 68--78.
https://doi.org/10.1037/0003-066X.55.1.68

Schulman, J., Wolski, F., Dhariwal, P., Radford, A., \& Klimov, O.
(2017). \emph{Proximal policy optimization algorithms}. arXiv preprint
arXiv:1707.06347.

Todd, G., Earle, S., Nasir, M. U., Green, M. C., \& Togelius, J. (2024).
PANGEA: Procedural artificial narrative using generative AI. In
\emph{Proceedings of the AAAI Conference on Artificial Intelligence and
Interactive Digital Entertainment (AIIDE 2024)}. AAAI Press.

Tondello, G. F., Wehbe, R. R., Diamond, L., Busch, M., Marczewski, A.,
\& Nacke, L. E. (2016). The gamification user types Hexad scale. In
\emph{Proceedings of the Annual Symposium on Computer-Human Interaction
in Play (CHI PLAY 2016)} (pp.~229--243). ACM.
https://doi.org/10.1145/2967934.2968082

Vornhagen, J. B., Thue, D., \& Göbel, S. (2022). miniPXI: Development
and validation of an eleven-item measure of player experience. In
\emph{Proceedings of the 2022 CHI Conference on Human Factors in
Computing Systems (CHI 2022)}, Article 329. ACM.
https://doi.org/10.1145/3491102.3502048

Wang, W., Wei, F., Dong, L., Bao, H., Yang, N., \& Zhou, M. (2020).
MiniLM: Deep self-attention distillation for task-agnostic compression
of pre-trained transformers. \emph{Advances in Neural Information
Processing Systems (NeurIPS), 33}, 5776--5788.

Yannakakis, G. N., \& Hallam, J. (2009). Real-time game adaptation for
optimizing player satisfaction. \emph{IEEE Transactions on Computational
Intelligence and AI in Games, 1}(2), 121--133.
https://doi.org/10.1109/TCIAIG.2009.2024533

Yee, N. (2006). Motivations for play in online games.
\emph{CyberPsychology \& Behavior, 9}(6), 772--775.
https://doi.org/10.1089/cpb.2006.9.772

\begin{center}\rule{0.5\linewidth}{0.5pt}\end{center}

\subsection{Appendices}\label{appendices}

\subsubsection{Appendix A: GEQ Core
Module}\label{appendix-a-geq-core-module}

The Game Experience Questionnaire Core Module as administered in this
study. Reproduced from IJsselsteijn et al.~(2013) for research purposes.

Full instrument file:
\texttt{E:\textbackslash{}Projects\textbackslash{}fyp\textbackslash{}v2\textbackslash{}evaluation\textbackslash{}questionnaires\textbackslash{}geq\_core.md}

\textbf{Scale:} 0 = Not at all, 1 = Slightly, 2 = Moderately, 3 =
Fairly, 4 = Very much

\emph{Please answer each question based on your experience during the
game session you just played.}

\begin{longtable}[]{@{}lll@{}}
\toprule\noalign{}
\# & Item & Scale \\
\midrule\noalign{}
\endhead
\bottomrule\noalign{}
\endlastfoot
1 & I felt content & 0--4 \\
2 & I felt skillful & 0--4 \\
3 & I was absorbed in the game & 0--4 \\
4 & I felt happy & 0--4 \\
5 & This game gave me a bad mood & 0--4 (R) \\
6 & I thought it was fun & 0--4 \\
7 & I was fully occupied with the game & 0--4 \\
8 & I felt annoyed & 0--4 (R) \\
9 & I found the game satisfying & 0--4 \\
10 & I felt frustrated & 0--4 (R) \\
11 & It felt like a whole new world & 0--4 \\
12 & I lost track of time & 0--4 \\
13 & I felt bored & 0--4 (R) \\
14 & I felt good & 0--4 \\
15 & I was good at this game & 0--4 \\
16 & I felt tired & 0--4 (R) \\
17 & I felt imaginative & 0--4 \\
18 & I felt that I could explore things & 0--4 \\
19 & I felt like I had to do things quickly & 0--4 \\
20 & I had fun with others & 0--4 \\
21 & I felt challenged & 0--4 \\
22 & I felt that time went by quickly & 0--4 \\
23 & I felt pressured & 0--4 (R) \\
24 & I felt competent & 0--4 \\
25 & I felt engrossed & 0--4 \\
26 & I felt bothered by aspects of the game & 0--4 (R) \\
27 & I was deeply concentrated in the game & 0--4 \\
28 & I forgot everything around me & 0--4 \\
29 & I felt pleased & 0--4 \\
30 & I felt completely focused on the game & 0--4 \\
31 & I enjoyed the game & 0--4 \\
32 & I felt imaginative & 0--4 \\
33 & I felt adventurous & 0--4 \\
\end{longtable}

\emph{(R) = reverse-scored}

\textbf{Subscale scoring:}

\begin{longtable}[]{@{}
  >{\raggedright\arraybackslash}p{(\columnwidth - 4\tabcolsep) * \real{0.3333}}
  >{\raggedright\arraybackslash}p{(\columnwidth - 4\tabcolsep) * \real{0.3333}}
  >{\raggedright\arraybackslash}p{(\columnwidth - 4\tabcolsep) * \real{0.3333}}@{}}
\toprule\noalign{}
\begin{minipage}[b]{\linewidth}\raggedright
Subscale
\end{minipage} & \begin{minipage}[b]{\linewidth}\raggedright
Items
\end{minipage} & \begin{minipage}[b]{\linewidth}\raggedright
Notes
\end{minipage} \\
\midrule\noalign{}
\endhead
\bottomrule\noalign{}
\endlastfoot
Flow (Primary DV, H1) & 5(R), 13(R), 25, 28, 31 & Mean of 5 items after
reverse-scoring \\
Immersion (Secondary DV, H2) & 3, 12, 18, 19, 27, 30 & Mean of 6
items \\
Positive Affect & 1, 4, 8(R), 20, 29 & Mean of 5 items \\
Negative Affect & 10, 16, 23, 26 & Mean of 4 items \\
Tension & 7, 22, 25 & Mean of 3 items \\
Competence & 2, 15, 17, 21, 24 & Mean of 5 items \\
\end{longtable}

\begin{center}\rule{0.5\linewidth}{0.5pt}\end{center}

\subsubsection{Appendix B: miniPXI
Questionnaire}\label{appendix-b-minipxi-questionnaire}

Full instrument file:
\texttt{E:\textbackslash{}Projects\textbackslash{}fyp\textbackslash{}v2\textbackslash{}evaluation\textbackslash{}questionnaires\textbackslash{}minipxi.md}

Citation: Vornhagen, J. B., Thue, D., \& Göbel, S. (2022). miniPXI:
Development and Validation of an Eleven-Item Measure of Player
Experience. \emph{Proc. ACM CHI} 2022. DOI:
https://doi.org/10.1145/3491102.3502048

\textbf{Scale:} 1 = Strongly Disagree \ldots{} 7 = Strongly Agree

\emph{The following statements describe how you felt while playing.
Please rate your agreement.}

\begin{longtable}[]{@{}ll@{}}
\toprule\noalign{}
\# & Item \\
\midrule\noalign{}
\endhead
\bottomrule\noalign{}
\endlastfoot
1 & I felt in control of my character/avatar \\
2 & I felt a sense of mastery when completing challenges \\
3 & I found the game world interesting to explore \\
4 & I was curious about what would happen next in the game \\
5 & I felt like the game narrative pulled me in \\
6 & I had a sense of presence in the game world \\
7 & I felt free to make meaningful choices \\
8 & The game felt meaningful to me \\
9 & I had positive emotions while playing \\
10 & I felt excited during the game \\
11 & Overall, I enjoyed playing this game \\
\end{longtable}

\textbf{Subscale mapping:}

\begin{longtable}[]{@{}ll@{}}
\toprule\noalign{}
Construct & Items \\
\midrule\noalign{}
\endhead
\bottomrule\noalign{}
\endlastfoot
Mastery & 1, 2 \\
Curiosity & 3, 4 \\
Immersion & 5, 6 \\
Autonomy & 7, 8 \\
Positive Affect & 9, 10, 11 \\
\end{longtable}

Total score: mean of all 11 items (range 1--7).

\begin{center}\rule{0.5\linewidth}{0.5pt}\end{center}

\subsubsection{Appendix C: NPS-3 Scale}\label{appendix-c-nps-3-scale}

Full instrument file:
\texttt{E:\textbackslash{}Projects\textbackslash{}fyp\textbackslash{}v2\textbackslash{}evaluation\textbackslash{}questionnaires\textbackslash{}personalization\_scale.md}

\emph{Custom Narrative Personalization Scale developed for this study.}

\textbf{Scale:} 1 = Strongly Disagree \ldots{} 7 = Strongly Agree

\emph{Please rate your agreement with each statement about the game's
story and narrative.}

\begin{longtable}[]{@{}ll@{}}
\toprule\noalign{}
\# & Item \\
\midrule\noalign{}
\endhead
\bottomrule\noalign{}
\endlastfoot
1 & The story felt tailored to how I was playing. \\
2 & The narrative responded to my choices and actions. \\
3 & The game seemed to understand what kind of player I am. \\
\end{longtable}

\textbf{Scoring:} NPS-3 Score = Mean of items 1--3 (range 1--7). Higher
scores indicate greater perceived personalization.

Cronbach's alpha will be reported; if α \textless{} 0.70, items will be
reported individually and the composite NPS-3 score treated as advisory
only.

\begin{center}\rule{0.5\linewidth}{0.5pt}\end{center}

\subsubsection{Appendix D: Interview
Guide}\label{appendix-d-interview-guide}

Full instrument file:
\texttt{E:\textbackslash{}Projects\textbackslash{}fyp\textbackslash{}v2\textbackslash{}evaluation\textbackslash{}questionnaires\textbackslash{}interview\_guide.md}

\textbf{Study:} AI-Driven Dynamic Narrative Personalization in Games
\textbf{Duration:} \textasciitilde20 minutes \textbf{Analysis method:}
Reflexive Thematic Analysis (Braun \& Clarke, 2022)

\textbf{Before the interview:} Obtain signed consent; remind participant
that there are no right or wrong answers and they can stop at any time;
start audio recording with consent.

\textbf{Section 1 --- General Experience (5 min)} 1. Can you walk me
through what you remember most about the game experience? 2. How would
you describe the story and the way it unfolded for you? 3. Were there
any moments that stood out to you, positively or negatively?

\textbf{Section 2 --- Narrative and Personalization (8 min)} 4. Did the
narrative feel connected to the way you were playing? Can you give an
example? 5. Were there moments where the story surprised you --- or felt
unexpectedly relevant? 6. Did any NPCs or dialogue feel particularly
fitting or particularly out of place? 7. \emph{(Experimental only)} Did
you notice the story adapting to you? How did that feel? 8.
\emph{(Control only)} How well did the story match your play style?

\textbf{Section 3 --- Flow and Engagement (5 min)} 9. Were there moments
where you felt in ``the zone'' --- neither too stressed nor too bored?
10. Were there moments where the game felt too easy, too hard, or just
right? What changed? 11. How did the story contribute to those moments
of engagement?

\textbf{Section 4 --- Closing (2 min)} 12. Is there anything about the
game or narrative you would like to add that the questionnaires might
not have captured? 13. Any suggestions for improving the experience?

\textbf{Probes:} ``Can you say more about that?'' / ``What do you mean
when you say {[}X{]}?'' / ``When did that happen in the game?'' / ``How
did that make you feel?''

\textbf{Post-interview memo:} After each interview, write a 3--5
sentence memo summarizing dominant themes, surprising responses, and
initial analytic observations.

\begin{center}\rule{0.5\linewidth}{0.5pt}\end{center}

\subsubsection{Appendix E: API
Endpoints}\label{appendix-e-api-endpoints}

The FastAPI backend exposes the following REST and WebSocket endpoints:

\begin{longtable}[]{@{}
  >{\raggedright\arraybackslash}p{(\columnwidth - 8\tabcolsep) * \real{0.2000}}
  >{\raggedright\arraybackslash}p{(\columnwidth - 8\tabcolsep) * \real{0.2000}}
  >{\raggedright\arraybackslash}p{(\columnwidth - 8\tabcolsep) * \real{0.2000}}
  >{\raggedright\arraybackslash}p{(\columnwidth - 8\tabcolsep) * \real{0.2000}}
  >{\raggedright\arraybackslash}p{(\columnwidth - 8\tabcolsep) * \real{0.2000}}@{}}
\toprule\noalign{}
\begin{minipage}[b]{\linewidth}\raggedright
Method
\end{minipage} & \begin{minipage}[b]{\linewidth}\raggedright
Path
\end{minipage} & \begin{minipage}[b]{\linewidth}\raggedright
Description
\end{minipage} & \begin{minipage}[b]{\linewidth}\raggedright
Request Body
\end{minipage} & \begin{minipage}[b]{\linewidth}\raggedright
Response
\end{minipage} \\
\midrule\noalign{}
\endhead
\bottomrule\noalign{}
\endlastfoot
POST & /telemetry/batch & Submit a telemetry window & TelemetryBatch
(JSON) & PlayerModel (JSON) \\
GET & /narrative/generate & Retrieve generated narrative for session &
session\_id (query param) & NarrativeContent (JSON) \\
GET & /player/model/\{player\_id\} & Get current player model & --- &
PlayerModel (JSON) \\
GET & /player/hexad/\{player\_id\} & Get current Hexad profile & --- &
HexadProfile (JSON) \\
DELETE & /session/\{session\_id\} & Clear session data & --- & 204 No
Content \\
WebSocket & /ws/telemetry/\{session\_id\} & Low-latency telemetry stream
& TelemetryBatch (JSON frames) & PlayerModel (JSON frames) \\
\end{longtable}

\textbf{TelemetryBatch (Request body --- POST /telemetry/batch):} All 23
fields as defined in
\texttt{E:\textbackslash{}Projects\textbackslash{}fyp\textbackslash{}v2\textbackslash{}backend\textbackslash{}models\textbackslash{}telemetry.py}.
Required fields: player\_id, session\_id, window\_start, window\_end.
All game metric fields default to 0.

\textbf{NarrativeContent (Response body):}

\begin{Shaded}
\begin{Highlighting}[]
\FunctionTok{\{}
  \DataTypeTok{"action\_taken"}\FunctionTok{:} \StringTok{"PROVIDE\_GUIDANCE"}\FunctionTok{,}
  \DataTypeTok{"content\_type"}\FunctionTok{:} \StringTok{"dialogue"}\FunctionTok{,}
  \DataTypeTok{"content"}\FunctionTok{:} \StringTok{"The path below is treacherous. If you still carry the torch from the entry hall, keep it lit."}\FunctionTok{,}
  \DataTypeTok{"speaker"}\FunctionTok{:} \StringTok{"Captain Aldric"}\FunctionTok{,}
  \DataTypeTok{"emotional\_tone"}\FunctionTok{:} \StringTok{"tense"}\FunctionTok{,}
  \DataTypeTok{"lore\_refs"}\FunctionTok{:} \OtherTok{[]}\FunctionTok{,}
  \DataTypeTok{"fallback"}\FunctionTok{:} \KeywordTok{false}
\FunctionTok{\}}
\end{Highlighting}
\end{Shaded}

\textbf{Error handling:} All endpoints return standard HTTP status
codes. 422 Unprocessable Entity is returned for malformed request bodies
(Pydantic validation). 503 Service Unavailable is returned if Ollama is
unreachable; in this case the narrative endpoint returns a fallback
NarrativeContent with \texttt{"fallback":\ true}.

\begin{center}\rule{0.5\linewidth}{0.5pt}\end{center}

\subsubsection{Appendix F: MDP
Specification}\label{appendix-f-mdp-specification}

\textbf{Formal MDP tuple:} (S, A, T, R, γ)

\textbf{State space S:} Box(7,) ∈ {[}0, 1{]}\^{}7

\begin{longtable}[]{@{}llll@{}}
\toprule\noalign{}
Dimension & Symbol & Range & Description \\
\midrule\noalign{}
\endhead
\bottomrule\noalign{}
\endlastfoot
s\_0 & flow\_score & \{0.1, 0.2, 0.3, 1.0\} & Encoded current Flow
state \\
s\_1 & csr\_norm & {[}0, 1{]} & challenge\_skill\_ratio / 3.0 \\
s\_2 & h\_explorer & {[}0, 1{]} & Hexad Explorer score \\
s\_3 & h\_socializer & {[}0, 1{]} & Hexad Socializer score \\
s\_4 & h\_achiever & {[}0, 1{]} & Hexad Achiever score \\
s\_5 & h\_disruptor & {[}0, 1{]} & Hexad Disruptor score \\
s\_6 & t\_progress & {[}0, 1{]} & session\_elapsed / 1800 \\
\end{longtable}

\textbf{Action space A:} Discrete(8)

\begin{longtable}[]{@{}ll@{}}
\toprule\noalign{}
Index & Action \\
\midrule\noalign{}
\endhead
\bottomrule\noalign{}
\endlastfoot
0 & LOWER\_STAKES \\
1 & RAISE\_STAKES \\
2 & ADD\_MYSTERY \\
3 & ADD\_HUMOR \\
4 & PROVIDE\_GUIDANCE \\
5 & INCREASE\_URGENCY \\
6 & LORE\_REWARD \\
7 & NO\_CHANGE \\
\end{longtable}

\textbf{Transition function T:} For helpful actions (see
\_HELPFUL\_ACTIONS in \texttt{backend/adaptation/rl\_env.py}): - P(FLOW
\textbar{} helpful action, non-FLOW state) = 0.75 - P(current state
\textbar{} helpful action, non-FLOW state) = 0.25 - P(FLOW \textbar{}
NO\_CHANGE, FLOW) = 0.90 - P(BOREDOM \textbar{} NO\_CHANGE, FLOW) = 0.10
- P(current state \textbar{} unhelpful action, non-FLOW state) = 1.0

\textbf{Reward function R(s, a, s'):}

R = R\_state(s') + R\_step + R\_streak(s', steps) + R\_transition(s, s')
+ R\_repetition(a, last\_3)

Where: - R\_state: FLOW = +1.0, BOREDOM = -0.3, ANXIETY = -0.5, APATHY =
-0.8 - R\_step = -0.05 (per step) - R\_streak = +0.20 if s' == FLOW AND
steps\_in\_state \textgreater{} 2 - R\_transition = +0.30 if s ∈
\{ANXIETY, BOREDOM\} AND s' == FLOW - R\_repetition = -0.15 if
all(last\_3\_actions identical)

\textbf{Discount factor γ:} 0.95

\textbf{Episode length:} max\_steps = 360 (one step per 5-second
telemetry window, corresponding to 30 minutes)

\textbf{Training:} PPO (Stable-Baselines3 v2.x), 200,000 timesteps, 4
parallel environments. See Table 9 for full hyperparameter
specification. Model checkpoint:
\texttt{./data/ppo\_narrative\_agent.zip}

\textbf{Cold-start policy (session\_elapsed \textless{} 120s):}

\begin{longtable}[]{@{}ll@{}}
\toprule\noalign{}
Flow State & Action \\
\midrule\noalign{}
\endhead
\bottomrule\noalign{}
\endlastfoot
ANXIETY & PROVIDE\_GUIDANCE \\
BOREDOM & INCREASE\_URGENCY \\
APATHY & ADD\_MYSTERY \\
FLOW & NO\_CHANGE \\
\end{longtable}

Implementation:
\texttt{E:\textbackslash{}Projects\textbackslash{}fyp\textbackslash{}v2\textbackslash{}backend\textbackslash{}adaptation\textbackslash{}agent.py}

\begin{center}\rule{0.5\linewidth}{0.5pt}\end{center}

\emph{End of thesis.}
